\chapter{经典自守形式}\label{ch:03}

\begin{center}
E. Kowalski
\end{center}

\section{引言}\label{c3:sec:1}

一旦涉足自守形式及其相伴的 L 函数，我们便进入了一片全新的领域。这里的口号是：描述 Dirichlet 特征的 $\GL(1)$ 理论（即 $\mathbb{Q}$ 上的理论）的 $\GL(2)$ 类比。至于更一般的群与一般基域 $K$ 上的相应工作，将在后面几章中介绍。

若要更完整地了解这套优美的理论，可参看 Iwaniec 的书（[I1]、[I2]）、Miyake 的书 [Mi] 或 Shimura 的书 [Sh]，Serre 的书 [Se1] 第 VII 章，以及 Bump 的书 [Bu] 第 I 章。

\section{自守形式的三个动机}\label{c3:sec:2}

我们先给出引入自守形式的三个理由。其中两个显然与前两章提到的问题相关，第三个则具有十分自然的代数几何内涵。自守形式的一个奇妙之处在于：人们还能找出许多表面互不相干的动机；此处所举的例子并未涉及与 Galois 表示的联系、同余、分拆等重要课题。

\subsection{Theta 函数}\label{c3:sec:2.1}

Riemann zeta 函数函数方程的证明（它是第 2 章相应证明的特例）要用到基本的 theta 函数
\begin{equation}\tag{2.1}\label{c3:eq:2.1}
\theta(z)=\sum_{n\in\mathbb{Z}}e\bigl(n^{2}z/2\bigr),
\end{equation}
其中 $z=iy$，$y>0$。事实上，它在到处出现的 Poincaré 上半平面 $\mathbf{H}=\{z\in\mathbb{C}\mid y>0\}$ 上定义了一个全纯函数。Poisson 求和公式（连同 $x\mapsto e^{-\pi x^{2}}$ 对 Fourier 变换自对偶这一事实）经解析延拓给出函数方程
\[
\theta\Bigl(-\frac{1}{z}\Bigr)=(-iz)^{1/2}\theta(z).
\]

其中 $(-iz)^{1/2}$ 取该函数在 $\mathbf{H}$ 上使 $iy\mapsto\sqrt{y}$ 的那一支。显然 $\theta(z)$ 还是 $1$ 周期的，故对由 $z\mapsto z+1$ 与 $z\mapsto-1/z$ 生成的群中任一映射 $z\mapsto\gamma z$，$f(\gamma z)$ 都有某种变换公式。这个群 $\Gamma$ 同构于按
\begin{equation}\tag{2.2}\label{c3:eq:2.2}
\begin{pmatrix}a&b\\c&d\end{pmatrix}\cdot z=\frac{az+b}{cz+d}
\end{equation}
作用于 $\mathbf{H}$ 的 $\mathrm{PSL}(2,\mathbb{Z})$。

同样可推出，对任意 $k\ge1$，$\theta(2z)^{k}$ 满足类似的变换公式。由于
\[
\theta(2z)^{k}=\sum_{n\ge0}r_{k}(n)e(nz),
\]
其中 $r_{k}(n)$ 是把 $n$ 表为 $k$ 个整数平方和的方式数，古典学界围绕这类函数投入大量兴趣也就不足为奇。诸如
\begin{equation}\tag{2.3}\label{c3:eq:2.3}
\begin{aligned}
r_{2}(n)&=4\sum_{d\mid n}\chi_{4}(d),\\
r_{4}(n)&=8\bigl(2+(-1)^{n}\bigr)\sum_{\substack{d\mid n\\ d\ \text{odd}}}d,\\
r_{6}(n)&=16\sum_{d\mid n}d^{2}\chi_{4}\Bigl(\frac{n}{d}\Bigr)-4\sum_{d\mid n}d^{2}\chi_{4}(d)
\end{aligned}
\end{equation}
这样的恒等式，可用（或多或少）初等或特设的办法推出；但模形式理论能解释这类公式的一般形状，能说明 $k$ 大时为何没有如此“初等”的公式，还能给出渐近关系，表明 $r_{k}(n)$ 的大小约为 $n^{k-1}$（$k\ge5$）。远为深刻的结果属于 Kloosterman（$k=4$ 个变量的一般二次型）与 Iwaniec（$k=3$ 个变量）（见 [I1, 5-3, 11-4]），例如它们给出球面上整点的等分布结果。

\subsection{行列式方程的解数}\label{c3:sec:2.2}

研究 $\zeta(s)$ 或其他 L 函数在临界带内的阶时，一条自然的解析途径是考察矩
\[
I_{k}(T)=\int_{-T}^{T}\bigl|\zeta(1/2+it)\bigr|^{k}\,dt,
\]
因为若能证明对一切 $k\ge1$ 都有
\begin{equation}\tag{2.4}\label{c3:eq:2.4}
I_{k}(T)\ll_{k,\varepsilon}T^{1+\varepsilon}\qquad(T\ge1),
\end{equation}
Lindelöf 猜想 $\zeta(1/2+it)\ll_{\varepsilon}(1+|t|)^{\varepsilon}$ 便容易推出。注意，即便对单个 $k\ge1$，(2.4) 也只是在平均意义下验证 Lindelöf 猜想，且 $k$ 越大结论越强，因而在解析数论中确有许多应用。

对 $k=2$，$I_{2}(T)$ 的渐近求值相当简单而古典（例如见 [Ti, 7.3]）；但 $k=4$ 的情形已颇具挑战，而对任何 $k>4$，(2.4) 均未知。$k=4$ 时，标准技巧把 (2.4) 归结为估计指数和
\[
\Bigl(\sum_{n\le X}n^{-it}\Bigr)^{4}\qquad(X\sim\sqrt{T})
\]
的四阶矩；把该和展开，得
\[
S=\sum_{a,b,c,d\le X}\Bigl(\frac{bc}{ad}\Bigr)^{it}.
\]
于是人们自然会按行列式 $h=ad-bc$ 的值把和拆开，即
\[
S=\sum_{|h|\le X^{2}}\ \sum_{ad-bc=h}\Bigl(1-\frac{h}{ad}\Bigr)^{it}.
\]
思路如下：若 $h=0$，便得方程 $ad=bc$，它易于求解，给出“主项”。若 $h\ne0$ 但“很小”，则指数 $(1-h/ad)^{it}$ 振荡不剧烈，也会带来重要贡献；而 $h$“很大”时，振荡应引起大量抵消。为计入小的 $h$ 的贡献，需要精确地数出方程
\[
ad-bc=h
\]
在 $a,b,c,d\le X$ 下的解数。用“通常的”调和分析即可做到，\footnote{实际上，这很快就牵涉到 Kloosterman 和；而无论用何种方式估计它们，背后都潜伏着自守形式。}但若注意到 $\mathrm{SL}(2,\mathbb{Z})$ 作用在解集上（当没有大小限制时），并对该作用施以某种调和分析（即某种形式的 Poisson 求和公式），则可得到好得多的结果。这里，我们从经典的 $\mathbb{Z}^{n}$ 离散于 $\mathbb{R}^{n}$ 的情形（及其在欧氏空间格点计数中的应用），走向 $\mathrm{SL}(2,\mathbb{Z})$ 离散于 $\mathrm{SL}(2,\mathbb{R})$ 的情形与相应的格点问题。然而后者的性质完全不同：$\mathrm{SL}(2,\mathbb{R})$ 非交换，整个几何是双曲的负曲率型（例如一个区域的绝大部分面积都积聚在其边界附近）。关于这些方面的更多内容，见 [I2]。

\subsection{椭圆曲线的不变量}\label{c3:sec:2.3}

考虑标准环面 $\mathbf{T}=\mathbb{R}^{2}/\mathbb{Z}^{2}$；它是二维光滑 Lie 群。问：$\mathbf{T}$ 上能赋哪些相容的复 Lie 群结构？这个问题的答案自然引向研究 $\mathbb{R}^{2}=\mathbb{C}$ 中的格在幺模变换下的等价类。

的确，若 $\Lambda\subset\mathbb{C}$ 是格，则商 $\mathbb{C}/\Lambda$ 微分同胚于 $\mathbf{T}$，并从 $\mathbb{C}$ 继承相容的复结构与群结构。反之，一切这样的结构都必由此得到。

由于 $\mathbb{C}$ 作为群的所有复自同构都是线性的 $z\mapsto\alpha z$，可见两个复环面 $\mathbb{C}/\Lambda_{1}$ 与 $\mathbb{C}/\Lambda_{2}$ 作为 Riemann 曲面同构，当且仅当 $\Lambda=\alpha\Lambda'$ 对某个 $\alpha\in\mathbb{C}^{\times}$ 成立。今设 $\Lambda=\omega_{1}\mathbb{Z}\oplus\omega_{2}\mathbb{Z}$，$\Lambda'=\omega'_{1}\mathbb{Z}\oplus\omega'_{2}\mathbb{Z}$，则 $\Lambda=\Lambda'$ 当且仅当两组基元由幺模线性变换相联系：
\[
\omega'_{1}=a\omega_{1}+b\omega_{2},\qquad
\omega'_{2}=c\omega_{1}+d\omega_{2},
\]
其中
\[
\gamma=\begin{pmatrix}a&b\\c&d\end{pmatrix}\in\mathrm{SL}(2,\mathbb{Z}),\quad\text{即 }ad-bc=1.
\]
在相似意义下，$\Lambda\simeq\mathbb{Z}\oplus\tau\mathbb{Z}$，$\tau=\omega_{1}/\omega_{2}$（相应地 $\Lambda'\simeq\mathbb{Z}\oplus\tau'\mathbb{Z}$，$\tau'=\omega'_{1}/\omega'_{2}$），于是方程化为
\[
\tau'=\gamma\cdot\tau=\frac{a\tau+b}{c\tau+d},
\]
最后我们还可以交换 $\omega_{1}$ 与 $\omega_{2}$，即把 $\tau$ 换成 $1/\tau$：两种选择中恰有一种落在 $\mathbf{H}$ 内。

于是，复环面的集合可等同于 $\mathbf{H}$ 对 $\mathrm{SL}(2,\mathbb{Z})$ 这一作用的商；这类复环面的各种不变量，便表现为在该作用下不变的函数 $f\colon\mathbf{H}\to\mathbb{C}$。传统上，以两个 $\mathbb{R}$ 线性无关的周期 $\omega_{1},\omega_{2}$ 为周期的 $\mathbb{C}$ 上的亚纯函数称为椭圆函数；许多椭圆函数自然地以如下方式出现：它们给出关于某权重 $k$ 齐次的函数 $f(z;\omega_{1},\omega_{2})$，即
\[
f(\lambda z;\lambda\omega_{1},\lambda\omega_{2})=\lambda^{-k}f(z;\omega_{1},\omega_{2}),
\]
从而
\[
f(u;\omega_{1},\omega_{2})=\omega_{2}^{-k}g(u/\omega_{2},\omega_{1}/\omega_{2}),
\]
其中 $g$ 定义在 $\mathbb{C}\times(\mathbf{H}\cup\overline{\mathbf{H}})$ 上，限制在 $\mathbb{C}\times\mathbf{H}$ 上时满足
\[
g(z;\gamma\tau)=(c\tau+d)^{k}g(z;\tau),\qquad\gamma\in\mathrm{SL}(2,\mathbb{Z}).
\]

\par\medskip\noindent\textbf{例 2.1}\quad 格 $\Lambda$ 的 Weierstrass $\wp$ 函数为
\[
\wp(z;\omega_{1},\omega_{2})=\frac{1}{z^{2}}+\sum_{\substack{\omega\in\Lambda\\ \omega\ne0}}\Bigl(\frac{1}{(z-\omega)^{2}}-\frac{1}{\omega^{2}}\Bigr).
\]

并且它是 $\Lambda$ 不变的（这并不显然：求导可见 $\wp'$ 确为 $\Lambda$ 不变，再积分即得）。显然，
\[
\wp(\lambda z;\lambda\omega_{1},\lambda\omega_{2})=\lambda^{-2}\wp(z;\omega_{1},\omega_{2}).
\]
此外
\[
\wp(z;\Lambda)=z^{-2}+\sum_{m\ge1}(m+1)z^{m}G_{m+2}(\Lambda),
\]
其中
\[
G_{k}(\Lambda)=\sum_{\omega\ne0}\frac{1}{\omega^{k}},
\]
而 $G_{k}(\lambda\Lambda)=\lambda^{-k}G_{k}(\Lambda)$ 给出这类模函数的又一例。椭圆函数这一主题同样十分古典，深受 18、19 世纪数学家的喜爱。

\begin{annote}
$\wp$ 的 Laurent 展开这里写作 $\sum_{m\ge1}(m+1)z^{m}G_{m+2}(\Lambda)$，看似与常见的 $\sum_{j\ge1}(2j+1)z^{2j}G_{2j+2}$ 不同，实则一致：$k$ 为奇数时 $G_{k}(\Lambda)=\sum_{\omega\ne0}\omega^{-k}=0$（$\omega$ 与 $-\omega$ 的项相消），故所有奇数 $m$ 的项自动消失，剩下的正是 $m=2j$ 的项。抄录时若把 $z^{m}$ 改回 $z^{2m}$ 而不改求和指标，就会得到错误的系数。
\par\vspace{0.15em}{\raggedleft 【来源：原书 §2.3（A）；Lang《Elliptic Functions》第 1--2 章（A）】\par}
\end{annote}

\section{定义与例子}\label{c3:sec:3}

下面我们简要综述所谓自守形式与 L 函数的经典理论的主要定义与结果。这一理论在历史上属于 20 世纪初，对我们最重要的名字是 Eisenstein、Poincaré、Hecke、Ramanujan、Petersson 以及（较晚的）Maass 与 Selberg。我们将同时谈论全纯模形式与非全纯（Maass）形式，尽管后者出现得晚得多：这恰好强烈地促使人们发展更具群论色彩的方法，以统一当时的工作。

为简单起见，我们大多只关注 $\mathrm{SL}(2,\mathbb{Z})$ 的同余子群，这是算术上最重要的情形。

回忆 $\mathbf{H}$ 是 Poincaré 上半平面。它是单连通的 Riemann 曲面；赋予 Riemann 度量
\[
ds^{2}=\frac{dx^{2}+dy^{2}}{y^{2}}
\]
后，它便是双曲平面（常负曲率 $-1$）的一个模型。

行列式为正的可逆矩阵构成的群 $\mathrm{GL}(2,\mathbb{R})^{+}$ 按 (2.2) 中的分式线性变换作用于 $\mathbf{H}$。纯量矩阵作用平凡，故人们常把注意力限于 $\mathrm{SL}(2,\mathbb{R})$ 或商群 $\mathrm{PGL}(2,\mathbb{R})^{+}=\mathrm{PSL}(2,\mathbb{R})$。后者是 $\mathbf{H}$ 作为 Riemann 曲面的自同构群，也是 $\mathbf{H}$ 作为 Riemann 流形的保向等距群。全部等距群是半直积
\[
\mathrm{PSL}(2,\mathbb{R})\rtimes\mathbb{Z}/2\mathbb{Z},
\]
其中反向等距的旁集的代表是关于虚轴的对称 $z\mapsto-\bar z$。

与 Riemann 度量相伴的是体积元 $d\mu(z)=y^{-2}dx\,dy$ 与 Laplace 算子 $\Delta=-y^{2}(\partial_{x}^{2}+\partial_{y}^{2})$，二者都在等距变换下不变。

对照 $\mathbb{R}$ 上周期函数或椭圆函数的研究，人们考虑离散子群\footnote{上下文总能清楚地区分子群与伽马函数。} $\Gamma<\mathrm{SL}(2,\mathbb{R})$，进而考虑 $\mathbf{H}$ 上关于 $\Gamma$“周期”的函数，即商空间 $\Gamma\backslash\mathbf{H}$ 上的函数，或微分形式，等等。

一个相当大的差别在于：这些商空间的形状千差万别。拓扑上，几乎所有（可定向）曲面都这样出现；而亏格 $>1$ 的 Riemann 曲面全体（自 Riemann 以来就知道）在复同构意义下依赖于 $3g-3$ 个复“参数”或“模”。

从算术角度看，子群 $\mathrm{SL}(2,\mathbb{Z})$ 引人注目：它确实离散地作用于 $\mathrm{SL}(2,\mathbb{R})$，因而 $\mathrm{SL}(2,\mathbb{Z})$ 的每个子群也如此。一般地，同余子群 $\Gamma<\mathrm{SL}(2,\mathbb{Z})$ 指满足 $\Gamma\supset\Gamma(q)=\{\gamma\in\mathrm{SL}(2,\mathbb{Z})\mid\gamma\equiv1\ (\mathrm{mod}\ q)\}$（对某个 $q$）者。同余群中，Hecke 同余群
\[
\Gamma_{0}(q)=\Bigl\{\gamma=\begin{pmatrix}a&b\\c&d\end{pmatrix}\in\mathrm{SL}(2,\mathbb{Z})\ \Big|\ q\mid c\Bigr\}
\]
（对任意整数 $q\ge1$）将特别重要。数 $q$ 称为该群的级。

$\Gamma_{0}(q)$ 在 $\Gamma_{0}(1)=\mathrm{SL}(2,\mathbb{Z})$ 中的指数容易算出：
\[
[\Gamma_{0}(1):\Gamma_{0}(q)]=q\prod_{p\mid q}\Bigl(1+\frac{1}{q}\Bigr).
\]

\begin{annote}
此处右端的分母印作 $q$，应为 $p$：$\Gamma_{0}(q)$ 的指数是 $q\prod_{p\mid q}(1+1/p)$。取 $q=6$，正确值为 $6\cdot\frac32\cdot\frac43=12$，而印式给出 $6(1+1/6)^{2}=49/6$，甚至不是整数，足以见其为笔误。译文照录原书。
\par\vspace{0.15em}{\raggedleft 【来源：原书 §3（A）；Iwaniec《Topics in Classical Automorphic Forms》第 2 章（A）】\par}
\end{annote}

在“交换”情形，$\mathbb{R}/\mathbb{Z}$ 是紧的；但商 $\mathrm{SL}(2,\mathbb{Z})\backslash\mathbf{H}$ 并非如此，这就引出了尖点这一重要概念。

我们将如下几何地描述尖点：对任意离散子群 $\Gamma<\mathrm{SL}(2,\mathbb{R})$，可用基本域来想象商空间 $\Gamma\backslash\mathbf{H}$，即 $\mathbf{H}$ 的一个开子集 $F$，使得
\begin{itemize}[leftmargin=2em]
\item $F$ 中任何两点都不 $\Gamma$ 等价；
\item 任何 $\Gamma$ 轨道 $\Gamma z$ 都与 $F$ 的闭包 $\overline{F}$ 相交。
\end{itemize}
在 $\mathrm{SL}(2,\mathbb{Z})$ 的情形，可取那个以无穷远点为顶点的熟悉三角形
\[
F_{1}=\{z=x+iy\in\mathbf{H}\mid-1/2<x<1/2,\ \text{且}\ |z|>1\}
\]
（例如见 [Se1, VII-1.2]）。

于是，对 $\mathrm{SL}(2,\mathbb{Z})$ 的子群 $\Gamma$（这正是我们将要处理的情形），可取
\[
F_{\Gamma}=\bigcup_{g\in\Gamma\backslash\mathbf{H}}g\cdot F_{1}.
\]

\begin{annote}
下标印作 $g\in\Gamma\backslash\mathbf{H}$，字面上讲不通：$\Gamma\backslash\mathbf{H}$ 的元素是轨道而非群元素，无法再作用。此处应理解为 $g$ 取遍 $\Gamma$ 在 $\mathrm{SL}(2,\mathbb{Z})$ 中的一组陪集代表，即 $g\in\Gamma\backslash\mathrm{SL}(2,\mathbb{Z})$；这样 $\bigcup_{g}g\cdot F_{1}$ 才是 $\Gamma$ 的基本域。译文照录原书，读者按陪集代表理解即可。
\par\vspace{0.15em}{\raggedleft 【来源：原书 §3（A）；Serre《Cours d'arithmétique》第 VII 章（A）】\par}
\end{annote}
Gauss--Bonnet 定理或直接积分表明 $F_{1}$ 的双曲体积有限（等于 $\pi/3$）；但由于 $in\in F_{1}$ 对 $n>1$ 成立，它并不紧。对 $\Gamma_{0}(q)$ 情形相同，只是体积要乘以 $\Gamma_{0}(q)$ 在 $\mathrm{SL}(2,\mathbb{Z})=\Gamma_{0}(1)$ 中的指数。

紧性的缺失当然出现在基本域达到“边界”$\overline{\mathbf{R}}=\mathbf{R}\cup\infty$ 的那些点处。注意 $\mathrm{GL}(2,\mathbb{R})^{+}$ 也按同一公式 (2.2) 作用于边界。

若 $\Gamma$ 的余体积 $\operatorname{Vol}(\Gamma\backslash\mathbf{H})=\operatorname{Vol}(F)$ 有限，则尖点只有有限多个；因为围绕每个尖点都能造出彼此不交、体积为正常数的小邻域。没有尖点当且仅当 $\Gamma\backslash\mathbf{H}$ 紧。（这里我们不会真正讨论这类群；不过存在若干极重要的算术例子，它们基于四元数代数的单位群，并出现在其他章将要讨论的 Jacquet--Langlands 对应中。）用群论的话说，群 $\Gamma$ 的任一尖点 $\mathfrak{a}$ 都是某个 $\gamma\in\Gamma$ 的唯一不动点（这样的元素称为抛物元）。这样的 $\gamma$ 在 $\mathbf{H}$ 上的几何作用是沿“环绕尖点”的极限圆作平移：对 $\mathfrak{a}=\infty$，
\begin{equation}\tag{3.1}\label{c3:eq:3.1}
\gamma=\begin{pmatrix}1&1\\0&1\end{pmatrix}\quad\text{作用为 }z\mapsto z+1,
\end{equation}
而极限圆是水平直线 $\operatorname{Im}(z)=y_{0}$。

由传递性，对任一尖点 $\mathfrak{a}$，存在伸缩矩阵 $\sigma_{\mathfrak{a}}\in\mathrm{SL}(2,\mathbb{R})$，使得
\[
\begin{cases}
\sigma_{\mathfrak{a}}\infty=\mathfrak{a},\\[2pt]
\sigma_{\mathfrak{a}}^{-1}\gamma_{\mathfrak{a}}\sigma_{\mathfrak{a}}=\begin{pmatrix}1&1\\0&1\end{pmatrix}.
\end{cases}
\]
借助共轭，关于尖点的一切概念都可化归为 $\mathfrak{a}=\infty$ 的情形。这更便于想象（与双曲平面的其他模型相比，Poincaré 上半平面有一个被区别出来的尖点，这是很方便的特性）。这里 $\gamma\in\Gamma_{\mathfrak{a}}<\Gamma$ 是尖点在 $\Gamma$ 中的稳定化子的一位生成元，该稳定化子是无限循环群（至少投影到 $\mathrm{PSL}_{2}$ 后如此）。

\par\medskip\noindent\textbf{定义}\quad 设 $f\colon\mathbf{H}\to\mathbb{C}$ 是 $\mathbf{H}$ 上的 $1$ 周期光滑函数（$f(z+1)=f(z)$），
\[
f(z)=\sum_{n\in\mathbb{Z}}a_{f}(n;y)e(nx)
\]
为其 Fourier 展开。称
\begin{enumerate}[label=(\roman*),leftmargin=2.4em]
\item $f$ 在 $\infty$ 处具有\emph{适度增长}，若存在 $A>0$ 使
\[
f(z)\ll|y|^{A}\qquad(y\ge1);
\]
\item $f$ 在 $\infty$ 处是\emph{尖点的}，若 $a_{0}(y)=0$。
\end{enumerate}

除 $\mathrm{SL}(2,\mathbb{R})$ 对 $\mathbb{C}$ 上函数的自然右作用 $(f\mid\gamma)(z)=f(\gamma z)$ 外，对任意整数 $k$ 还定义作用
\[
f\mapsto f\mid_{k}\gamma,\qquad\text{其中 }(f\mid_{k}\gamma)(z)=(cz+d)^{-k}f(\gamma z).
\]
可以验证确实有 $f\mid_{k}\gamma\mid_{k}\gamma'=f\mid_{k}\gamma\gamma'$。此外，矩阵 (3.1) 的作用仍为 $(f\mid_{k}\gamma)(z)=f(z+1)$。

至此可以定义模形式了。

\par\medskip\noindent\textbf{定义}\quad 设 $k\in\mathbb{Z}$，$\lambda\in\mathbb{C}$，$q\ge1$ 为整数，$\chi$ 为模 $q$ 的 Dirichlet 特征。$\mathbf{H}$ 上的光滑函数 $f\colon\mathbf{H}\to\mathbb{C}$ 称为权 $k$、级 $q$、nebentypus 为 $\chi$ 的\emph{模形式}，相应地，称为特征值 $\lambda$、级 $q$、nebentypus 为 $\chi$ 的\emph{自守函数}，若 $f$ 在 $\mathbf{H}$ 上全纯且满足
\[
f\mid_{k}\gamma=\chi(a)f,\qquad\gamma=\begin{pmatrix}a&b\\c&d\end{pmatrix}\in\Gamma_{0}(q),
\]
相应地，$f$ 满足 $\Delta f=\lambda f$ 且
\[
f\mid\gamma=\chi(a)f,\qquad\gamma\in\Gamma_{0}(q).
\]

\par\medskip\noindent\textbf{注 3.1}\quad 取 $\gamma=-1$，得关系 $f=f\mid_{k}\gamma=\chi(-1)(-1)^{k}f$，故非零全纯模形式只能存在于特征满足相容条件 $\chi(-1)=(-1)^{k}$ 之时；下文默认如此。对自守函数同理，必须有 $\chi(-1)=1$。

\begin{annote}
“nebentypus” 是德文旧称，中文常译“旁特征”。注 3.1 的相容条件来自 $-1\in\Gamma_{0}(q)$：权 $k$ 的作用在 $-1$ 上给出因子 $(-1)^{k}$，而特征给出 $\chi(-1)$，二者必须相等，否则 $f=-f$，形式只能为零。Maass 形式用权 $0$ 的作用，故只要求 $\chi(-1)=1$。
\par\vspace{0.15em}{\raggedleft 【来源：原书 §3（A）；Iwaniec《Topics in Classical Automorphic Forms》第 2--3 章（A）】\par}
\end{annote}

也可以借助修正的微分算子，对权 $k\ne0$ 定义非全纯形式。全纯形式与 Maass 形式都可经由 $\mathrm{GL}(2,\mathbb{R})$（算术情形下则是阿代尔群 $\mathrm{GL}(2,\mathbb{A})$）的表示论，被最有说服力地纳入同一框架。

利用上面的定义，还可以对尖点加更多正则性条件。

\par\medskip\noindent\textbf{定义}\quad 设 $f$ 为模形式（相应地，自守函数）。称 $f$ 在尖点处\emph{全纯}（相应地，为 \emph{Maass 形式}），当且仅当它在所有尖点处具有适度增长：即对任一 $\mathfrak{a}$，$1$ 周期函数 $f_{\mathfrak{a}}=f\mid_{k}\sigma_{\mathfrak{a}}$ 在无穷远处具有适度增长。以 $M_{k}(q,\chi)$（相应地 $M^{\lambda}(q,\chi)$）记权 $k$、级 $q$、特征 $\chi$ 的模形式（相应地 Maass 形式）构成的向量空间。

称 $f$ 为\emph{尖点形式}（相应地 \emph{Maass 尖点形式}），若它在所有尖点处都是尖点的；并以 $S_{k}(q,\chi)\subset M_{k}(q,\chi)$（相应地 $S^{\lambda}(q,\chi)$）记尖点形式（相应地 Maass 尖点形式）的子空间。

\par\medskip\noindent\textbf{注 3.2}\quad 还可给出其他等价表述。特别地，对 $\Gamma_{0}(q)$，自守函数 $f$ 是 Maass 尖点形式，当且仅当 $f$ 非常数（即 $\lambda\ne0$）且属于 $L^{2}(\Gamma_{0}(q)\backslash\mathbf{H})$（关于双曲测度）。这等价于说 $0$ 是 $\Gamma_{0}(q)$ 唯一的剩余特征值（例如见 [I2, 11.2]）。

对全纯形式，无穷远处的 Fourier 展开形如
\begin{equation}\tag{3.2}\label{c3:eq:3.2}
f(z)=\sum_{n\in\mathbb{Z}}a_{f}(n)e(nz),\qquad a_{f}(n)\in\mathbb{C},
\end{equation}
而 $f$ 全纯意味着 $n<0$ 时 $a_{n}=0$，因为 $|e(nz)|=\exp(-2\pi ny)$ 对 $n<0$ 不是多项式有界的。

\par\medskip\noindent\textbf{注 3.3}\quad 令 $Y_{0}(q)=\Gamma_{0}(q)\backslash\mathbf{H}$。这是一个非紧的 Riemann 曲面（该作用在 $\mathbf{H}$ 中有一些不动点，它们都 $\mathrm{SL}(2,\mathbb{Z})$ 等价于 $i$（由 $z\mapsto-1/z$ 固定）或 $\exp(i\pi/3)$（由 $z\mapsto(z-1)/z$ 固定）。然而在这些点处，适当的坐标仍给出复结构；细节见 [Sh, 第 1 章]）。

此外，添上有限多个尖点可将 $Y_{0}(q)$ 紧化，得到紧 Riemann 曲面 $X_{0}(q)$。例如在 $\infty$ 处用 $q=e(z)$ 作坐标卡（与 (3.2) 比较）。当 $q=1$ 时，借助 $j$ 不变量 $j(z)$ 得 $X_{0}(1)\simeq\mathbf{P}(1)$。

由 Riemann 曲面的一般理论可知，$X_{0}(q)$ 实际上是代数曲线。但还有更多：这条代数曲线实际定义在 $\mathbb{Q}$ 上，并且在 $\mathbb{Z}$ 上有很好的模型（“坏纤维”也了解得相当清楚）。这大多基于把 $X_{0}(q)$ 解释为\emph{模空间}：它“分类”对 $(E,H)$，其中 $E$ 是椭圆曲线，$H$ 是 $E$ 上 $q$ 阶循环子群。点 $z\in\mathbf{H}$ 在 $X_{0}(q)$ 中的像对应于 $(\mathbb{C}/(\mathbb{Z}\oplus z\mathbb{Z}),\langle1/q\rangle)$。这些都是极重要的内容，将在后面几章讨论，也是模形式许多深刻算术性质的关键；它与我们关于模形式的第三个动机相关。

\begin{annote}
读本章时须分清两种 $q$：一是 $\Gamma_{0}(q)$ 的级（正整数），二是尖点 $\infty$ 处的局部坐标 $q=e(z)$。后者在注 3.3、例 4.1 与 6.2 节前后反复出现，且只在 $|q|<1$ 时有意义；“$q$-展开”一语即指按 $e(z)$ 展开。混用这两个 $q$ 是初读时最常见的误会。
\par\vspace{0.15em}{\raggedleft 【来源：本库术语与译法（I）；Iwaniec《Topics in Classical Automorphic Forms》第 2 章（A）】\par}
\end{annote}

这类联系的一个简单例子是重要的同构
\[
\begin{cases}
S_{2}(q)\longrightarrow\Gamma(X_{0}(q),\Omega_{1}),\\
f\longmapsto f(z)\,dz,
\end{cases}
\]
它把级 $q$（nebentypus 平凡）的权 $2$ 尖点形式空间与模曲线 $X_{0}(q)$（而非 $Y_{0}(q)$）上的全纯 $1$ 形式空间联系起来。这相当容易证明：权 $2$ 的作用恰好是用来说明 $f(z)dz$ 在 $\Gamma_{0}(q)$ 下不变所需的，从而它下降到 $X_{0}(q)$ 上的 $1$ 形式。此外，对 $\infty$ 处的局部坐标 $q=e(z)$ 有 $dq/q=2\pi i\,dz$，故在尖点 $\infty$ 处消没的条件等价于 $f(z)dz$ 在 $\infty$ 处全纯，其他尖点同理。

本注最后一段表明（因 $X_{0}(q)$ 紧，整体微分构成的空间 $\Gamma(X_{0}(q),\Omega_{1})$ 有限维），$\dim_{\mathbb{C}}S_{2}(q)<+\infty$。事实上，$\dim M_{k}(q,\chi)$（相应地 $\dim M^{\lambda}(q,\chi)$）有限是普遍事实。\footnote{这在远为一般的框架下也成立。}

对全纯形式，可用 Cauchy 定理给出初等论证（例如见 [Se1, VII-3]）；但对 $k>1$，包括 $\dim M_{k}(q,\chi)$ 精确公式在内的更深结果，则如 [Sh, 2.6] 那样由 Riemann--Roch 定理推出：例如由上述，$\dim S_{2}(q)$ 就是 $X_{0}(q)$ 的亏格，而利用 $\Gamma_{0}(q)<\Gamma_{0}(1)$ 诱导的自然（分歧）覆叠 $X_{0}(q)\to X_{0}(1)=\mathbf{P}(1)$，可以显式算出它。

除小幅波动外，维数随级与权一起增长：$\dim M_{k}(q)$ 的大小（粗略地）为 $qk/12$，故只要权或级不太小，模形式就很多（例如 $X_{0}(q)$ 的亏格在 $q>49$ 时 $>1$，或在 $q$ 为素数 $>19$ 时 $>1$）。

被排除的 $k=1$ 情形\footnote{$k\le0$ 时没有模形式。}确实相当不同：计算其维数具有重大算术意义，且在很大程度上仍是公开问题（见 [Se2] 与 [Du]）。从谱的观点看，权 $k=1$ 对应连续谱中的特征值，故用迹公式之类的调和分析工具极难把它挑出来；任何合理的试验函数也会探测到与它接近的特征值的（很大）贡献。

对 Maass 形式，$\dim M^{\lambda}(q,\chi)$ 的有限性借助 Hilbert 空间中算子的谱理论证明。与全纯情形相反，这只是一个定性结论，该空间的维数没有任何已知公式。这是可以理解的，因为谱理论还蕴含：除（至多！）可数个趋于 $+\infty$ 的特征值 $\lambda>0$ 外，$S^{\lambda}(q,\chi)=0$（而 $M^{\lambda}/S^{\lambda}$ 易于描述）。这些特征值完全神秘，只知少数特例（$\lambda=1/4$ 是特例，也具有深刻的算术意义——这句话在此反复出现；另外几个可由实二次域的 Hecke 特征构造出相伴的 Maass 形式，这是 Maass 证明的；例子见 [Ge]）。

Selberg 发展他著名的迹公式，至少部分是为了处理 Maass 尖点形式的存在性问题。借助它，他证明了（例如见 [I2, 第 11 章]）：对同余子群，特别是对 $\Gamma_{0}(q)$，Weyl 律成立：
\[
\bigl|\{\lambda\mid0\le\lambda\le X\ \text{且}\ S^{\lambda}(q)\ne0\}\bigr|\sim\frac{\operatorname{Vol}(X_{0}(q))}{4\pi}\,X
\]
当 $X\to+\infty$，特征值按重数计。\footnote{如何精确地给出 Maass 尖点形式重数的上界，是解析数论中最难以捉摸的公开问题之一。}

\begin{annote}
原书此处写 $\operatorname{Vol}(X_{0}(q))$，而正文其他地方用的是 $\Gamma_{0}(q)\backslash\mathbf{H}$ 的体积。二者其实相等：紧化时添入的尖点只有有限多个，在双曲测度下是零测集，故按 $\Gamma_{0}(q)\backslash\mathbf{H}$ 计算不会出错。要紧的是别把 $X_{0}(q)$ 的亏格（例 3.6 中用到）与它的体积混为一谈。
\par\vspace{0.15em}{\raggedleft 【来源：原书 §3（A）；Iwaniec《Introduction to the Spectral Theory of Automorphic Forms》第 10--11 章（A）】\par}
\end{annote}

另一方面，现在人们倾向于认为（这与 Selberg 证明上述结果时的预期相反）：对“一般的”（非算术）群 $\Gamma$，这一论断会以极糟糕的方式失效，以致 $\Gamma$ 的 Maass 尖点形式只有有限多个（甚至可能没有）特征值。这一普遍看法的转变，归功于 Phillips 与 Sarnak 关于尖点形式在群形变下消失的优美理论 [PS]。

本节最后，我们给出各类模形式的一些具体例子。

\par\medskip\noindent\textbf{例 3.4}\quad 构造在群 $G$ 的某个作用下“周期”的函数，最自然的办法是对函数在群上求平均（参看 Poisson 求和公式的证明，即命题 2.1）。事实上，若函数已在子群 $G_{1}<G$ 下不变，则只需对 $G_{1}$ 在 $G$ 中的陪集求平均。由此引出 Poincaré 级数与 Eisenstein 级数。

设 $m\ge0$ 为整数。定义权 $k$ 的第 $m$ 个全纯 Poincaré 级数为
\[
P_{m}(z)=\sum_{\gamma\in\Gamma_{\infty}\backslash\Gamma_{0}(q)}\overline{\chi}(a)(cz+d)^{-k}e(m\gamma z).
\]
当 $k>2$ 时，该级数在紧集上绝对一致收敛到一个 $\mathbf{H}$ 上的全纯函数；由于上述理由，它是级 $q$ 的模形式。可以算出 $P_{m}$ 在各尖点处的展开，从而看出 $P_{m}\in M_{k}(q,\chi)$，且事实上当 $m\ne0$ 时 $P_{m}\in S_{k}(q,\chi)$（见下文）。当 $m=0$ 时，称之为 Eisenstein 级数。取 $q=1$，便得到经典的 Eisenstein 级数
\[
E_{k}(z)=\sum_{(a,b)=1}\frac{1}{(az+b)^{k}}.
\]
Eisenstein 级数与 Poincaré 级数也可在非全纯情形下定义，这时把 $k$ 换成复变量 $s$：
\begin{equation}\tag{3.3}\label{c3:eq:3.3}
P_{m}(z;s)=\sum_{\gamma\in\Gamma_{\infty}\backslash\Gamma_{0}(q)}\overline{\chi}(a)\operatorname{Im}(\gamma z)^{s}e(m\gamma z),\qquad E(z;s)=P_{0}(z;s).
\end{equation}

（注意对 $\gamma\in\mathrm{SL}(2,\mathbb{R})$ 有 $\operatorname{Im}(\gamma z)=|cz+d|^{-2}$。）它在 $\operatorname{Re}(s)>1$ 时收敛，并在 $\mathbf{H}$ 上定义实解析函数，满足
\[
P_{m}(\cdot;s)\mid\gamma=\overline{\chi}(d)P_{m}(\cdot;s)\qquad(\gamma\in\Gamma_{0}(q)),
\]
但当 $m\ne0$ 时它不是 Laplace 算子的特征函数；事实上，
\[
\Delta P_{m}(\cdot;s)=\bigl(s(1-s)-4\pi^{2}m^{2}\bigr)P_{m}(\cdot;s)+4\pi ms\,P_{m}(\cdot;s+1).
\]
另一方面，Eisenstein 级数 $E(z;s)$ 在 $\operatorname{Re}(s)>1$ 时确实是 $\Delta$ 的特征函数，特征值为 $\lambda=s(1-s)$，且容易验证它在各尖点处具有多项式增长。正如后面几章将要讨论的，非全纯 Eisenstein 级数即使在研究全纯形式时也起着非常重要的作用，例如通过它们在 Rankin--Selberg 方法中的出现（这一方法的简短介绍见 6.2 节）。

一般地，把任一特征值写成 $s(1-s)$ 的形式往往很方便（于是给定 $\lambda$ 有两个 $s$ 值），并进一步令 $s=1/2+it$，$t\in\mathbb{C}$。对尖点形式，$\lambda>0$ 转化为 $\lambda\in]0,1[$（若 $\lambda\le1/4$）或 $\lambda=1/2+it$，$t\in\mathbb{R}$。第一种情形下，该特征值称为\emph{异常的}。Selberg 提出了如下猜想：

\begin{annote}
原文这里的两处 $\lambda$ 均应为 $s$：由 $\lambda=s(1-s)$，当 $0<\lambda\le1/4$ 时两根均为实数，取 $s\in[1/2,1]$；当 $\lambda>1/4$ 时写作 $s=1/2+it$，$t\in\mathbb{R}$。若照字面读成“$\lambda=1/2+it$”，便与 $\lambda$ 为实数自相矛盾，可见是笔误。译文照录原书。
\par\vspace{0.15em}{\raggedleft 【来源：原书 §3（A）；Iwaniec《Introduction to the Spectral Theory of Automorphic Forms》第 11 章（A）】\par}
\end{annote}

\par\medskip\noindent\textbf{猜想 3.5}\quad 对任意同余子群 $\Gamma$，特别是 $\Gamma_{0}(q)$，不存在异常特征值；即 $\Delta$ 作用在 $L^{2}(\Gamma\backslash\mathbf{H})$ 上的第一个非零特征值满足 $\lambda_{1}\ge1/4$。

此外，Selberg 证明了 $\lambda_{1}\ge3/16$。这是一个算术命题，因为容易构造出具有任意多个异常特征值的非同余子群。在 2.2 节所提的那类问题中，$\Gamma_{0}(q)$ 的异常特征值所起的作用类似于 Dirichlet 特征的 Landau--Siegel 零点：关于 $q$ 的一致性（例如在数方程 $ad-bc=h$ 满足 $c\equiv0\ (\mathrm{mod}\ q)$、$a^{2}+b^{2}+c^{2}+d^{2}\le X$ 的解时）会受到存在“许多”异常特征值的影响（越接近 $0$，影响越坏）。因此 Selberg 的定理表明情况略好控制一些。这一猜想今天也被理解为 Ramanujan--Petersson 猜想的阿基米德类比（见下文 4.2 节）；事实上，Luo、Rudnick 与 Sarnak [LRS] 对 $3/16$ 界所作的重要改进，正是基于这一类比。

对一般 Poincaré 级数的谱分析，是 Kuznetsov--Bruggeman 公式的精髓；该公式在 Maass 形式对解析数论的应用中极为重要（见 [I2, 第 9 章]、[CP]）。

严格说来，我们定义 Poincaré 级数与 Eisenstein 级数时是相对于尖点 $\infty$ 的。对任一尖点都可定义类似的函数；而 Eisenstein 级数 Fourier 展开的计算表明
\[
M_{k}(q,\chi)=S_{k}(q,\chi)\oplus E_{k}(q,\chi),
\]
其中 $E_{k}(q,\chi)$ 是由（各尖点处的）Eisenstein 级数张成的向量空间。

\par\medskip\noindent\textbf{例 3.6}\quad 现在我们来看 theta 函数。记法如 [I1, 第 10 章]，讨论一个相当一般的情形。设 $A$ 是对称正定的 $(r\times r)$ 整系数矩阵，其对角元素全为偶数，$N\ge1$ 是使 $NA^{-1}$ 为整系数矩阵的任一整数。令
\[
A[x]={}^{t}xAx,\qquad x\in\mathbb{R}^{r},
\]
为相伴的二次型。设变量个数 $r$ 为偶数。\footnote{$r$ 为奇数时，理论将引向半整数权模形式；因篇幅与时间所限，我们不拟介绍。}令
\[
\Theta(z;A)=\sum_{m\in\mathbb{Z}^{r}}e\bigl(A[m]z/2\bigr).
\]
\par\medskip\noindent\textbf{命题 3.7}\quad 有 $\Theta(z;A)\in M_{r/2}(2N)$。

这用 Poisson 求和公式证明，自然推广了公式 (2.5)。

\begin{annote}
本章的带号公式止于 (2.4)，并无 (2.5)。所指数应即 2.1 节中 theta 函数的变换公式 $\theta(-1/z)=(-iz)^{1/2}\theta(z)$（原书未编号），高秩的 $\Theta(z;A)$ 正是它的推广。此处存疑，译文照录原书。
\par\vspace{0.15em}{\raggedleft 【来源：原书 §2.1、§3（A）】\par}
\end{annote}

例如，取 $A_{4}=\operatorname{diag}(2,2,2,2)$。可取 $N=2$，此时 $A_{4}[x]/2=x_{1}^{2}+x_{2}^{2}+x_{3}^{2}+x_{4}^{2}$，故 $\Theta(z;A_{4})=\theta(2z)^{4}$，其中 $\theta$ 是基本 theta 函数 (2.1)。于是 $\Theta(\cdot;A_{4})\in M_{2}(4)$。今 $X_{0}(4)$ 的亏格为零，故 $S_{2}(4)=0$；于是如上例所解释，$\Theta\in E_{2}(4)$。$\Gamma_{0}(4)$ 有三个尖点，即 $0$、$1/2$ 与 $\infty$；可以显式算出权 $2$ 的三个相应 Eisenstein 级数 $E_{0}$、$E_{1/2}$ 与 $E_{\infty}$ 的 Fourier 展开。核对前几个系数，即可定出等于 $\Theta(z;A_{4})$ 的那个线性组合：
\[
\Theta(z;A_{4})=\alpha E_{\infty}(z)+\beta E_{0}(z),
\]
其中 $\alpha,\beta$ 可显式写出；由此即得 (2.3)。

\par\medskip\noindent\textbf{例 3.8}\quad Ramanujan Delta 函数为
\[
\Delta(z)=e(z)\prod_{n\ge1}\bigl(1-e(nz)\bigr)^{24}.
\]
有 $\Delta\in S_{12}(1)$，且事实上该空间是一维的。利用 Eisenstein 级数 $E_{k}\in M_{k}(1)$，还有另一表达式：
\[
\Delta=(60E_{4})^{3}-27(140E_{6})^{2}.
\]
用椭圆曲线的语言说，若 $E\simeq\mathbb{C}/(\mathbb{Z}\oplus z\mathbb{Z})$，则 $\Delta(z)$（差一个非零常数）就是曲线 $E$ 的判别式。

\section{Fourier 展开与 L 级数}\label{c3:sec:4}

\subsection{定义}\label{c3:sec:4.1}

我们已经提到，模形式在 $\mathbf{H}$ 上特别是 $1$ 周期的，因此存在如下形式的 Fourier 展开

\begin{equation}\tag{4.1}\label{c3:eq:4.1}
f(z)=\sum_{n\ge0}a_{f}(n)e(nz)\qquad(\text{对全纯形式 }f),
\end{equation}
\begin{equation}\tag{4.2}\label{c3:eq:4.2}
f(x+iy)=\sum_{n\in\mathbb{Z}}a_{f}(n;y)e(nx)\qquad(\text{对非全纯形式}),
\end{equation}
用伸缩矩阵 $\sigma_{\mathfrak{a}}$ 作共轭后，对所考虑的群的每个尖点都有类似的展开。

在 (4.2) 的情形，若 $f$ 还满足 $\Delta f=\lambda f$，则对两边作用 $\Delta$（分离变量！）便得到 $a_{f}(n;y)$ 所满足的二阶常微分方程
\[
y^{2}w''+\bigl(\lambda-4\pi^{2}n^{2}y^{2}\bigr)w=0,
\]
它属于 Bessel 型。若 $n\ne0$，两个线性无关的解是
\[
w_{1}(y)=\sqrt{y}K_{s-1/2}(2\pi|n|y),\qquad
w_{2}(y)=\sqrt{y}I_{s-1/2}(2\pi|n|y)
\]
（这里写作 $\lambda=s(1-s)$，$s\in\mathbb{C}$），其中 $K_{s}$ 与 $I_{s}$ 表示标准的 Bessel 函数（例如见 [I2, 附录 B-4]）。$n=0$ 时，两个解是 $y^{s}$ 与 $y^{1-s}$（当 $s=1/2$ 时则为 $y^{1/2}$ 与 $y^{1/2}\log y$）。

重要的事实是 $I_{s}$ 与 $K_{s}$ 在无穷远处的渐近行为不同：$w_{1}(y)\sim\frac{\pi}{2}e^{-2\pi|n|y}$，而 $w_{2}(y)\sim\frac{1}{2\pi}e^{2\pi|n|y}$（这里没有笔误，只是各种极不一致的规范化留下的遗产）。因此对适度增长的 Maass 形式，第二个解在 $n\ne0$ 时不能出现；我们把 $f$ 的（在 $\infty$ 处的）Fourier 展开写成如下规范形式：
\begin{equation}\tag{4.3}\label{c3:eq:4.3}
f(z)=a_{f}(0)y^{s}+b_{f}(0)y^{1-s}+\sqrt{y}\sum_{n\ne0}a_{f}(n)K_{s-1/2}(2\pi|n|y)e(nx),
\end{equation}
其中 $a_{f}(n)\in\mathbb{C}$，$b_{f}(0)\in\mathbb{C}$。它表示一个在 $\infty$ 处为尖点的形式，当且仅当 $a_{f}(0)=b_{f}(0)=0$。

就 theta 函数而言，无穷远处的 Fourier 系数与整数用二次型表示的问题相关，显然具有算术意义。然而对更一般的形式，人们为何要关心 Fourier 系数、它们为何会有特殊性质，却一点也不显然。但事实证明确实如此；这一点至少可以部分地通过 L 函数显示出来，其方式与用 theta 函数研究 Hecke L 函数密切相关。

\subsection{例子；阶的大小}\label{c3:sec:4.2}

不过在朝这个方向走之前，我们先举出 Fourier 系数的一些显式计算。

\par\medskip\noindent\textbf{例 4.1}\quad 设 $k\ge3$，$m\ge0$ 为整数。Poincaré 级数 $P_{m}(z)$ 的第 $n$ 个 Fourier 系数（在 $\infty$ 处）对 $n\ge1$ 由下式给出：
\begin{equation}\tag{4.4}\label{c3:eq:4.4}
p(m,n)=\Bigl(\frac{m}{n}\Bigr)^{(k-1)/2}\Bigl\{\delta(m,n)+2\pi i^{-k}\sum_{q\mid c}c^{-1}S_{\overline{\chi}}(m,n;c)J_{k-1}\Bigl(\frac{4\pi\sqrt{mn}}{c}\Bigr)\Bigr\}
\end{equation}
（证明见例如 [I1, 3.2]），其中 $J_{k-1}$ 又是 Bessel 函数，$S_{\overline{\chi}}(m,n;c)$ 是扭 Kloosterman 和
\[
S_{\overline{\chi}}(m,n;c)=\sum_{x\ (\mathrm{mod}\ c)}^{*}\overline{\chi}(x)e\Bigl(\frac{mx+n\bar x}{c}\Bigr),
\]
求和取遍可逆元，$\bar x$ 表示 $x$ 模 $c$ 的逆。第 $0$ 个 Fourier 系数恰为零，其他尖点处的系数亦然，可见 $P_{m}$ 是尖点形式。

Bessel 函数也许显得古怪陌生，但有
\[
J_{k-1}(y)=O\bigl(y^{k-1}\bigr)\qquad(y\to0),
\]
（由幂级数展开即知）这说明上述级数收敛。此外，各种积分表示会揭示：$J_{k-1}(y)$ 实际上正是 Kloosterman 和的阿基米德类比（或者反过来）。这一点再次在阿代尔语言中解释得最好（见 [CP]）。

利用 $S_{k}(q,\chi)$ 上的 Petersson 内积
\begin{equation}\tag{4.5}\label{c3:eq:4.5}
\langle f,g\rangle=\int_{\Gamma_{0}(q)\backslash\mathbf{H}}f(z)\overline{g(z)}y^{k}\,d\mu(z)
\end{equation}
（被积函数是 $\Gamma_{0}(q)$ 不变的），可以证明
\[
\langle f,P_{h}\rangle=c_{k}h^{1-k}a_{f}(h),
\]
其中 $c_{k}>0$ 为（显式的）常数。由此可见 Poincaré 级数张成 $S_{k}(q,\chi)$。不过它们所满足的关系相当神秘。

\par\medskip\noindent\textbf{例 4.2}\quad 对 $\mathrm{SL}(2,\mathbb{Z})$，权为 $k$（偶）的 Eisenstein 级数在 $\infty$ 处的 Fourier 展开为
\begin{equation}\tag{4.6}\label{c3:eq:4.6}
E_{k}(z)=2\zeta(k)+\frac{2(2i\pi)^{k}}{(k-1)!}\sum_{n\ge1}\sigma_{k-1}(n)e(nz),
\end{equation}
其中
\[
\sigma_{k}(n)=\sum_{d\mid n}d^{k}.
\]
对 $\mathrm{SL}(2,\mathbb{Z})$ 的非全纯 Eisenstein 级数，有
\[
E(z,s)=y^{s}+\varphi(s)y^{1-s}+4\sqrt{y}\sum_{n\ge1}\eta_{s-1/2}(2\pi n)K_{s-1/2}(2\pi ny)\cos(2\pi nx),
\]

其中
\[
\eta_{u}(n)=\sum_{ab=n}\Bigl(\frac{a}{b}\Bigr)^{u},
\]
而 $\varphi(s)$ 是 $\mathrm{SL}(2,\mathbb{Z})$ 的散射矩阵，这里约化为单个函数：
\[
\varphi(s)=\sqrt{\pi}\frac{\Gamma(s-1/2)}{\Gamma(s)}\frac{\zeta(2s-1)}{\zeta(2s)}.
\]
该展开表明（借助 zeta 函数的性质），$E(z,s)$ 可解析延拓到整个 $\mathbb{C}$，并满足
\[
E(z,s)=\varphi(s)E(z,1-s),
\]
或由 $\zeta(s)$ 的函数方程，对称地写成
\begin{equation}\tag{4.7}\label{c3:eq:4.7}
\tilde{\varphi}(s)E(z,s)=\tilde{\varphi}(1-s)E(z,1-s),\qquad\text{其中 }\tilde{\varphi}(s)=\pi^{-s}\Gamma(s)\zeta(2s).
\end{equation}
这类结果在远为一般的框架下成立：Selberg 对 $\mathrm{SL}(2,\mathbb{R})$ 的有限余体积子群证明了它，Langlands 又将其推广到高秩群。Eisenstein 级数的解析延拓是自守形式的谱理论及其在解析数论中应用的关键之一（例如见 [I2, 第 6 章]）。

\par\medskip\noindent\textbf{例 4.3}\quad $\Delta$ 在 $\mathrm{SL}(2,\mathbb{Z})$ 的唯一尖点 $\infty$ 处的第 $n$ 个 Fourier 系数记为 $\tau(n)$；算术函数 $\tau$ 通常称为 Ramanujan 函数。

由于所考虑的模形式具有适度增长，容易证明其 Fourier 系数 $a_{f}(n)$ 作为 $n$ 的函数至多是多项式大小的：若 $f$ 全纯，则对任意 $x\in\mathbb{R}$、$y>0$ 有
\[
a_{f}(n)=\int_{x}^{x+1}f(x+t+iy)e(-nt)\,dt.
\]
对 Eisenstein 级数，精确的 Fourier 展开表明（忽略算术波动，见 (4.6)），$a_{f}(n)$ 的大小约为 $n^{k-1}$。

出于同样的理由，尖点形式的系数应当更小，因为 $f$ 在无穷远处消没。事实确实如此：对有界函数 $y^{k/2}f(z)$（相应地，非全纯情形下的 $f(z)$）应用 Parseval 公式，即可证明
\begin{equation}\tag{4.8}\label{c3:eq:4.8}
\sum_{n\le N}|a_{f}(n)|^{2}\ll N^{k}\qquad(N\ge1),
\end{equation}
（取 $y=N^{-1}$）；相应地，若 $f$ 非全纯，则 $\ll N$。

事实表明这已接近真相（作为平均结果），正如 Rankin--Selberg 方法所显示的；它还表明，平均说来至少有 $a_{f}(n)\ll n^{(k-1)/2}$（相应地，平均有 $a_{f}(n)\ll1$）。

这一点的个体形式首先由 Ramanujan 就 $\tau$ 函数猜想，且形式惊人地精确：
\begin{equation}\tag{4.9}\label{c3:eq:4.9}
|\tau(n)|\le d(n)n^{11/2}
\end{equation}
（$d(n)=\sigma_{0}(n)$ 是 $n$ 的除数个数）。

Petersson 把该猜想推广到其他模形式（形式不如前者精确）：应当有
\begin{equation}\tag{4.10}\label{c3:eq:4.10}
a_{f}(n)\ll_{\varepsilon}n^{(k-1)/2+\varepsilon},\qquad(n\ge1),
\end{equation}
对任意 $f\in S_{k}(q,\chi)$ 成立，隐含常数依赖于 $f$ 与 $\varepsilon>0$。

确实，除了上述平均意义上的证据，把“开平方”哲学施于 (4.4) 中 Kloosterman 和的和式，也可迅速为此期望提供依据；因为 Poincaré 级数张成 $S_{k}(q,\chi)$（$k>2$ 时），人们由此得到 (4.10) 的有力证据。然而沿这些线路始终未能找到证明；P. Deligne [De] 是通过对有限域上簇的 Riemann 假设的一个非凡应用（其中涉及模形式与 Galois 表示或“动机”的联系），才对权 $\ge2$ 的尖点形式证明了 Ramanujan--Petersson 猜想，即对 $\tau$ 函数得到 (4.10) 的精确形式，或更一般地
\[
|a_{f}(n)|\le d(n)n^{(k-1)/2}
\]
对本原形式（见下文第 5 节）的系数成立。权为 $1$ 时，相应的结果由 Deligne 与 Serre 证明。

不过也不应忽视经由 (4.4) 的途径；它的类比仍适用于例如 Maass 形式或半整数权形式，而在这些情形下相应猜想仍未解决（半整数权情形见 [I1, 5.3]，它对三变量二次型有非凡的应用）。

\subsection{模形式相伴随的 L 级数}\label{c3:sec:4.3}

设 $f$ 是权为 $k$ 的模形式或 Maass 形式，于是它在 $\infty$ 处有 Fourier 展开 (4.1) 或 (4.3)（也可在其他尖点处工作）。为研究系数的性质，考虑相伴的 L 级数（看来最好把“L 函数”一语留给带 Euler 乘积的 Dirichlet 级数）
\begin{equation}\tag{4.11}\label{c3:eq:4.11}
L(f,s)=\sum_{n\ge1}a_{f}(n)n^{-s}\qquad(\text{对全纯 }f),
\end{equation}
\begin{equation}\tag{4.12}\label{c3:eq:4.12}
L(f,s)=\sum_{n\ne0}a_{f}(n)|n|^{-s}\qquad(\text{对 Maass 形式}).
\end{equation}
Hecke（相应地，Maass\footnote{若 Maass 形式 $f$ 是奇的，即 $f(-\bar z)=-f(z)$，则该 L 级数消没；稍作修改仍然可行，我们不拟讨论。}）证明了这些 L 级数仍承载着前几章所考虑的 L 函数的大部分解析性质。首先注意：由 Fourier 系数的平凡界 $a_{f}(n)\ll n^{k}$（相应地 $a_{f}(n)\ll|n|$），该 L 级数在 $\operatorname{Re}(s)$ 充分大时绝对收敛。

令
\[
W=\begin{pmatrix}0&-1/\sqrt{q}\\ \sqrt{q}&0\end{pmatrix}\in\mathrm{SL}(2,\mathbb{R}),
\]
其作用为 $z\mapsto-1/(qz)$。这个 $W$ 正规化 $\Gamma_{0}(q)$（当 $q=1$ 时它属于 $\mathrm{SL}(2,\mathbb{Z})$），故简单计算表明：若 $f\in M_{k}(q,\chi)$，则 $g=f\mid_{k}W\in M_{k}(q,\overline{\chi})$（相应地，若 $f\in M^{\lambda}(q,\chi)$，则 $g=f\mid W\in M^{\lambda}(q,\overline{\chi})$）。

\par\medskip\noindent\textbf{命题 4.4}\quad 令
\[
\Lambda(f,s)=(2\pi)^{-s}\Gamma(s)L(f,s)
\]
（对全纯 $f$），
\begin{equation}\tag{4.13}\label{c3:eq:4.13}
\Lambda(f,s)=\pi^{-s}\Gamma\Bigl(\frac{s+it}{2}\Bigr)\Gamma\Bigl(\frac{s-it}{2}\Bigr)L(f,s)
\end{equation}
（对特征值 $\lambda=1/4+t^{2}$ 的 Maass 形式）。则 $\Lambda(f,s)$ 延拓为 $\mathbb{C}$ 上的亚纯函数。若 $f$ 是尖点形式，则它是整函数；否则它在 $s=1$ 与 $s=0$ 处有极点。此外还有
\[
\Lambda(f,s)=i^{k}q^{k/2-s}\Lambda(g,k-s),
\]
其中 $g=f\mid_{k}W$（对全纯 $f$）；或
\[
\Lambda(f,s)=q^{1/2-s}\Lambda(g,1-s),
\]
其中 $g=f\mid W$（对 Maass 形式 $f$）。

\begin{annote}
原书 (4.13) 左端印作 $L(f,s)$，与命题开头对 $\Lambda$ 的定义不符，应为 $\Lambda(f,s)$；否则该式只是把 $L$ 用自身表出，毫无内容。译文按定义改正，此处出注说明。另可注意两式的伽马因子不同：全纯情形是 $\Gamma(s)$，Maass 情形是 $\Gamma\bigl(\frac{s+it}{2}\bigr)\Gamma\bigl(\frac{s-it}{2}\bigr)$，这正是权 $k$ 与谱参数 $t$ 的差别在阿基米德位的体现。
\par\vspace{0.15em}{\raggedleft 【来源：原书 §4.3（A）；Iwaniec《Topics in Classical Automorphic Forms》第 5 章（A）】\par}
\end{annote}

\emph{证明}\quad 先考虑全纯情形。Maass 形式情形类似，只是须用 Bessel $K$ 函数的 Mellin 变换公式。我们还假设 $f$ 为尖点形式，以免处理常数项（其实也不难）。

证明的进行方式与证明 Hecke L 函数的解析延拓大致相同，只是以待研究的 $f$ 代替 theta 函数。对 $\operatorname{Re}(s)>1$ 有
\[
(2\pi)^{-s}\Gamma(s)n^{-s}=\int_{0}^{+\infty}e^{-2\pi ny}y^{s}\,\frac{dy}{y},
\]
故当 $\operatorname{Re}(s)$ 充分大时
\begin{equation}\tag{4.14}\label{c3:eq:4.14}
\Lambda(f,s)=\int_{0}^{\infty}f(iy)y^{s}\,\frac{dy}{y}.
\end{equation}

再把上述积分在 $\alpha>0$ 处拆开，得
\begin{align*}
\Lambda(f,s)&=\int_{0}^{\alpha}f(iy)y^{s}\,\frac{dy}{y}+\int_{\alpha}^{\infty}f(iy)y^{s}\,\frac{dy}{y}\\
&=\int_{1/(q\alpha)}^{+\infty}f\Bigl(\frac{i}{qu}\Bigr)(qu)^{-s}\,\frac{du}{u}+\int_{\alpha}^{\infty}f(iy)y^{s}\,\frac{dy}{y}\\
&=\int_{1/(q\alpha)}^{+\infty}q^{k/2}i^{k}u^{k}g(iu)(qu)^{-s}\,\frac{du}{u}+\int_{\alpha}^{\infty}f(iy)y^{s}\,\frac{dy}{y}\\
&=i^{k}q^{k/2-s}\int_{1/(q\alpha)}^{+\infty}g(iu)u^{k-s}\,\frac{du}{u}+\int_{\alpha}^{\infty}f(iy)y^{s}\,\frac{dy}{y},
\end{align*}
其中用到 $g(iu)=q^{-k/2}i^{-k}u^{-k}f(i/(qu))$。

这已经给出解析延拓；若取 $\alpha=q^{-1/2}$，便得到函数方程。\hfill$\square$

与头两章所见关于解析延拓与函数方程的各种结果相比，人们会注意到 gamma 因子不同（对 Maass 形式尤其如此），且函数方程把 $s$ 与 $k-s$ 相联系，从而临界带的中线位于 $s=k/2$。此外，函数方程把 $f$ 与 $g=f\mid_{k}W$ 联系起来。若 $q=1$，则 $W\in\mathrm{SL}(2,\mathbb{Z})$，故 $f=g$；但一般情形下没有理由认为这样的关系成立。

在解析数论中，常便于把系数重新规范化，令
\[
a_{f}(n)=n^{(k-1)/2}\lambda_{f}(n),
\]
并用
\[
\sum_{n\ge1}\lambda_{f}(n)n^{-s}
\]
代替 $L(f,s)$；于是这另一个级数的函数方程把 $s$ 与 $1-s$ 相联系。

\section{Hecke 算子及其应用}\label{c3:sec:5}

命题 4.4 给出了更多具有函数方程的 Dirichlet 级数的例子。然而似乎没有任何理由指望有 Euler 乘积。不过 Ramanujan 还猜想 $\tau$ 函数是乘性的（即当 $(m,n)=1$ 时 $\tau(mn)=\tau(m)\tau(n)$），更精确地，
\begin{equation}\tag{5.1}\label{c3:eq:5.1}
L(\Delta,s)=\sum_{n\ge1}\tau(n)n^{-s}=\prod_{p}\bigl(1-\tau(p)p^{-s}+p^{11-2s}\bigr)^{-1}.
\end{equation}

同样，$\theta(2z)^{2}$ 的 Fourier 系数是 $4r(n)$，其中 $r(n)$ 是 $\mathbb{Q}(i)$ 的 Dedekind zeta 函数的系数，于是
\begin{equation}\tag{5.2}\label{c3:eq:5.2}
\begin{aligned}
\sum_{n\ge1}r(n)n^{-s}&=\zeta(s)L(\chi_{4},s)=\prod_{p}(1-p^{-s})^{-1}(1-\chi_{4}(p)p^{-s})^{-1}\\
&=\prod_{p}\bigl(1-(1+\chi_{4}(p))p^{-s}+p^{-2s}\bigr)^{-1}.
\end{aligned}
\end{equation}
注意两种情形下分母关于 $p^{-s}$ 都是二次的（第二种情形下 $p=2$ 除外，此时 $\chi_{4}(2)=0$）。

Mordell 首先证明了 Ramanujan 的乘性猜想，但为建立一套融贯的理论、揭示全纯形式 Fourier 系数真正非凡的算术结构，贡献最大的是 Hecke。不过他未能对某些问题给出完全满意的答案；直到 Atkin 与 Lehner 发展了“新形式”理论，局面才真正明朗。Jacquet--Langlands 的阿代尔理论也为这些问题提供了许多启示。

Hecke 的想法是：通过在模形式空间上找到一个可对角化的作用代数，来获得“好的”算术模形式。特征函数于是继承该代数的大部分结构性质。对 Dirichlet 特征而言，这在“古典”意义下并无对应物，尽管事后可以构造出类比。

Hecke 算子有多种定义方式。我们举出几种：
\begin{itemize}[leftmargin=2em]
\item 可以对属于 $\Gamma$ 的可通约化子（commensurator）的某个 $\gamma$，由双陪集 $\Gamma\gamma\Gamma$ 生成一个抽象代数。对同余子群，该可通约化子大于 $\Gamma$。然后证明这个代数以适当方式作用于 $M_{k}(q,\chi)$，本质上是
\[
f\mid_{k}\Gamma\gamma\Gamma=\sum_{\alpha_{i}}f\mid_{k}\alpha_{i},
\]
其中 $\Gamma\gamma\Gamma=\bigcup_{i}\Gamma\alpha_{i}$ 是双陪集的分解。由局部--整体原则（中国剩余定理），可证明该代数由
\[
T_{p}=\Gamma\begin{pmatrix}p&0\\0&1\end{pmatrix}\Gamma\quad(p\ \text{为素数}),\qquad
R_{p}=\Gamma\begin{pmatrix}p&0\\0&p\end{pmatrix}\Gamma\quad(p\ \text{为素数})
\]
生成，且它们满足简单的乘性关系（见下文）。基于这一进路的详细研究见 [Sh]；在研究定义于不同子群上的 Hecke 算子之间的相互作用时，这一进路往往最为精确。
\item 若把模形式看作格空间上的函数，则可如下定义 Hecke 算子 $T_{p}$：
\[
(f\mid_{k}T_{p})(\Lambda)=\sum_{[\Lambda:\Lambda']=p}f(\Lambda'),
\]

其中求和取遍 $\Lambda$ 中指数为 $p$ 的所有子格（见 [Se1, VII-5]）。
\item 为了快速给出定义（虽然用它来证明算子的性质并不实用），可以直接给出它在 Fourier 系数上的作用；这里用到无穷远处的展开足以决定模形式这一事实。为简单起见，我们采用这一进路。
\end{itemize}

固定 $q$、权 $k$ 与特征 $\chi$。对素数 $p$，令
\[
T(p)=\frac{1}{p}\sum_{b\ (\mathrm{mod}\ p)}\begin{pmatrix}1&b\\0&p\end{pmatrix}+\chi(p)p^{k-1}\begin{pmatrix}p&0\\0&1\end{pmatrix}
\]
（在 $\mathrm{GL}(2,\mathbb{R})^{+}$ 的有理群环中）。当 $p\mid q$ 时注意第二项消没（即便对平凡特征，$\chi$ 也被视为模 $q$ 定义）。处理 Maass 形式时须采用不同的定义，即去掉因子 $p^{k-1}$；下文各公式作同样的改动。

对给定的权 $k$，我们让这些 $T(p)$ 以显然的方式（用 $\mid_k$ 作用或 $\mid$ 作用）作用于 $\mathbf{H}$ 上的 $1$ 周期函数，并记之为 $f\mid_{k}T(p)$ 或 $f\mid T(p)$。（注意 $T(p)$ 仍依赖于 $\chi$ 与 $q$。）

\par\medskip\noindent\textbf{引理 5.1}\quad 上述算子 $T(p)$ 作用于 $1$ 周期函数，诱导出 $M_{k}(q,\chi)$（相应地 $M^{\lambda}(q,\chi)$）上的自同态，并保持尖点形式。若 $a_{f}(n)$ 是 $f$ 的 Fourier 系数，则 $f\mid_{k}T(p)$ 的第 $n$ 个 Fourier 系数等于
\[
a_{f}(n/p)+\chi(p)p^{k-1}a_{f}(np),
\]
其中约定：当 $x\in\mathbb{Q}$ 不是整数时 $a_{f}(x)=0$。\footnote{如前，此公式针对全纯形式；对 Maass 形式须作改动（去掉 $p^{k-1}$）。}

本引理的第一部分当然才是关键论断；若不把我们所给的 $T(p)$ 的权宜定义与某个更内在的定义联系起来，就无法证明它。另一方面，第二部分相当简单（$\mathbb{Z}/p\mathbb{Z}$ 的特征的正交性）。

现在用乘性与归纳把定义推广到 $T(n)$，$n\ge1$：
\begin{align*}
T(nm)&=T(n)T(m)\qquad((n,m)=1),\\
T(p^{i+1})&=T(p)T(p^{i})+\chi(p)p^{k-1}T(p^{i-1}),\qquad(i\ge0),\\
T(1)&=1.
\end{align*}
简单的计算表明，$f\mid_{k}T(n)$ 的第 $m$ 个 Fourier 系数是
\begin{equation}\tag{5.3}\label{c3:eq:5.3}
\sum_{d\mid(n,m)}\chi(d)d^{k-1}a_{f}\Bigl(\frac{nm}{d^{2}}\Bigr);
\end{equation}
特别地，$f\mid_{k}T(n)$ 的第一个 Fourier 系数等于 $a_{f}(n)$，即 $f$ 的第 $n$ 个系数，而第 $0$ 个 Fourier 系数是 $\sigma_{k-1}(n)a_{f}(0)$。

\par\medskip\noindent\textbf{引理 5.2}\quad 由 Hecke 算子 $T(p)$ 或 $T(n)$ 生成的 $\operatorname{End}M_{k}(q,\chi)$ 的子代数是交换的，且有
\[
\sum_{n\ge1}T(n)n^{-s}=\prod_{p}\bigl(1-T(p)p^{-s}+\chi(p)p^{k-1-2s}\bigr)^{-1}.
\]
把这个 Euler 乘积的形状与 Ramanujan 对 $\tau$ 函数所猜想的 (5.1) 以及 (5.2) 中的形状作比较。事实上，我们有如下结果。

\par\medskip\noindent\textbf{推论 5.3}\quad 当 $\operatorname{Re}(s)$ 充分大时，Euler 乘积 (5.1) 成立。

\emph{证明}\quad 考虑权 $12$、级 $1$、$\chi=1$ 的 Hecke 算子。有 $\dim S_{12}(1)=1$，故 $S_{12}(1)=\Delta\mathbb{C}$。于是对任意 $n$ 有 $\Delta\mid_{12}T(n)=\lambda(n)\Delta$。由上述引理，我们同义反复地得到
\[
\sum_{n\ge1}\lambda(n)n^{-s}=\prod_{p}\bigl(1-\lambda(p)p^{-s}+p^{11-2s}\bigr)^{-1}.
\]
用 (5.3) 比较第一个 Fourier 系数，即得
\[
\tau(n)=\lambda(n)\tau(1)=\lambda(n),
\]
于是结论成立。\hfill$\square$

同理，对任何是所有 Hecke 算子 $T(n)$ 的特征函数的 $f$，都会有 Euler 乘积。但这样的 $f$ 存在吗？当权或级很大时，线性无关的全纯模形式很多，故推论的简单论证无法推广多远。然而 Hecke 证明了如下结果。

\par\medskip\noindent\textbf{引理 5.4}\quad 设 $n\ge1$ 与 $q$ 互素。则作用在 $S_{k}(q,\chi)$（相应地 $S^{\lambda}(q,\chi)$）上的算子 $T(n)$ 关于 Petersson 内积 (4.5)（相应地关于 $L^{2}(\Gamma_{0}(q)\backslash\mathbf{H})$ 中的内积）是正规的；事实上，$T(n)$ 的伴随算子是
\[
T(n)^{*}=\overline{\chi}(n)T(n).
\]
本引理的证明其实相当微妙。最重要的是，条件 $(n,q)=1$ 是必要的：当 $p\mid n$ 时，算子 $T(p)$ 通常并不正规。无论如何，Hecke 能推出：

\par\medskip\noindent\textbf{推论 5.5}\quad $S_{k}(q,\chi)$（相应地 $S^{\lambda}(q,\chi)$）有一组正交基，由所有满足 $(n,q)=1$ 的 Hecke 算子 $T(n)$ 的特征函数 $f$ 组成。这样的模形式称为 Hecke 形式，其 Hecke 特征值记为 $\lambda_{f}(n)$。

对 Hecke 形式 $f$，由方程
\[
f\mid_{k}T(n)=\lambda_{f}(n)f\qquad((n,q)=1)
\]
可以推出：当 $(n,q)=1$ 时 $a_{f}(n)=\lambda_{f}(n)a_{f}(1)$。若 $a_{f}(1)\ne0$，则由 Hecke 算子的 Euler 乘积推出，形式 $g=a_{f}(1)^{-1}f$ 满足
\begin{align*}
L_{q}(g,s)&=\sum_{(n,q)=1}a_{g}(n)n^{-s}=\sum_{(n,q)=1}\lambda_{f}(n)n^{-s}\\
&=\prod_{p\nmid q}\bigl(1-\lambda_{f}(p)p^{-s}+\chi(p)p^{k-1-2s}\bigr)^{-1}.
\end{align*}
还剩两个困难：可能 $a_{f}(1)=0$，此时上述推理失效；而我们更希望对 $L(f,s)$ 本身得到涉及所有素数的 Euler 乘积。

Hecke 能够证明，在某些情形下存在由所有 $T(n)$ 的特征函数构成的基；此时该基中形式的 L 函数会有 Euler 乘积（把第一个系数规范化之后）。特别地，若特征 $\chi$ 模 $q$ 本原，则确实如此；但这排除了 $\chi=1$ 而 $q>1$ 这一重要情形。Hecke 的想法是证明 Hecke 代数满足“重数一”：具有给定特征值 $\lambda(n)$（$(n,q)=1$）的模形式空间至多一维。若此事成立，则由于 $p\nmid q$ 的 Hecke 算子 $T(p)$ 作用在这个公共特征空间上（它们与所有其他 Hecke 算子交换），便推出对任意 $p$，$f\mid_{k}T(p)$ 是 $f$ 的倍数，理论即可如前推进。然而重数一在一般情形下失效。

Atkin 与 Lehner [AL] 指出，如何通过引入本原特征概念的一个类比来修正 Hecke 的理论。这很合理，因为我们知道只有本原 Dirichlet 特征才满足漂亮的函数方程。

该理论的出发点是如下把模形式诱导到更高级的办法：设 $q\ge1$，$\chi$ 为模 $q$ 的特征，$a>1$。设 $f\in S_{k}(q,\chi)$。则对任一整数 $d\mid a$，函数
\[
g_{d}(z)=f(dz)
\]
是 $S_{k}(aq,\chi')$ 中的元素，其中 $\chi'$ 是由 $\chi$ 诱导到模 $aq$ 的特征；且若 $f$ 是 Hecke 形式，则 $g_{d}$ 也是 Hecke 形式，并对 $(n,aq)=1$ 有相同的特征值。这一结果极易验证。特别地，$f$ 与 $g_{d}$ 提供了重数一失效的例子。（另注意
\[
g_{d}(z)=\sum_{n\equiv0\ (\mathrm{mod}\ d)}a_{f}(n/d)e(nz),
\]
故 $d>1$ 时 $g_{d}$ 的第一个 Fourier 系数为 $0$。）这个例子引出如下定义：

\par\medskip\noindent\textbf{定义}\quad 设 $q\ge1$，$\chi$ 为模 $q$ 的特征，$k$ 为整数。设 $q'\mid q$ 是 $\chi$ 的导子，$\chi'$ 是诱导出 $\chi$ 的本原特征。$S_{k}(q,\chi)$ 的\emph{旧空间}，记作 $S_{k}(q,\chi)^{b}$，是由所有形如 $f(az)$ 的函数张成的子空间，其中
\begin{enumerate}[label=(\roman*),leftmargin=2.6em]
\item $f\in S_{k}(q'',\chi')$，对某个满足 $q'\mid q''\mid q$、$q''\ne q$ 的 $q''$；
\item $a\mid q/q''$。
\end{enumerate}

\emph{新空间} $S_{k}(q,\chi)^{*}$ 是旧空间在 $S_{k}(q,\chi)$ 中关于 Petersson 内积的正交补。

\par\medskip\noindent\textbf{例 5.6}\quad (i) 若 $q=1$ 或 $\chi$ 本原，则 $S_{k}(q,\chi)=S_{k}(q,\chi)^{*}$。
(ii) 若 $k<12$ 且 $q=p$ 为素数，则由 $S_{k}(1)=0$ 推知 $S_{k}(p)^{*}=S_{k}(p)$。若 $k\ge12$，则
\[
S_{k}(p)=S_{k}(1)\oplus pS_{k}(1)\oplus S_{k}(p)^{*},
\]
其中 $pS_{k}(1)$ 记形如 $f(pz)$（$f\in S_{k}(1)$）的尖点形式构成的空间。一般地，可用 Möbius 反演计算新空间的维数。

容易看出，所有 $(n,q)=1$ 的 Hecke 算子 $T(n)$ 都作用在旧空间与新空间上。

\par\medskip\noindent\textbf{定理 5.7（Atkin--Lehner）}\quad 重数一原则对 $S_{k}(q,\chi)^{*}$ 成立：即若 $(\lambda(n))$ 是对 $(n,q)=1$ 定义的任一复数序列，则满足
\[
f\mid_{k}T(n)=\lambda(n)f\qquad((n,q)=1)
\]
的 $f\in S_{k}(q,\chi)^{*}$ 构成的空间至多一维。

$S_{k}(q,\chi)^{*}$ 有唯一的正交基，由本原形式（也称新形式）组成，即满足 $a_{f}(1)=1$ 且对任意 $n\ge1$ 有 $f\mid_{k}T(n)=a_{f}(n)f$ 的形式 $f$。

第一部分是关键论断，第二部分由上述梗概即可推出。

\par\medskip\noindent\textbf{推论 5.8}\quad 设 $f$ 为本原形式。则 L 函数 $L(f,s)$ 有 Euler 乘积
\begin{equation}\tag{5.4}\label{c3:eq:5.4}
L(f,s)=\prod_{p}\bigl(1-a_{f}(p)p^{-s}+\chi(p)p^{k-1-2s}\bigr)^{-1}.
\end{equation}
此外存在模长为 $1$ 的 $\eta(f)\in\mathbb{C}$，使得 $f\mid_{k}W=\eta(f)\tilde f\in S_{k}(q,\overline{\chi})^{*}$，其中
\[
\tilde f(z)=\sum_{n\ge1}\overline{a_{f}(n)}e(nz),
\]
故 $L(f,s)$ 满足函数方程
\[
\Lambda(f,s)=\varepsilon(f)q^{k/2-s}\Lambda(\tilde f,k-s),
\]
其中 $\varepsilon(f)=i^{k}\eta(f)$。

这个函数方程与 (3.3) 中本原 Dirichlet 特征的情形完全类似。$\varepsilon(f)$ 是一个非常微妙的完全不变量，一般没有简单公式。若 $q$ 无平方因子，则可把 $\eta(f)$ 与 $f$ 的第 $q$ 个 Fourier 系数联系起来。例如 $q=1$ 时，$f\mid_{k}W=f$，函数方程的符号为 $i^{k}$。

由此还可推出：事实上对任意级为 $q$ 的 Hecke 形式 $f$，都存在唯一的本原形式 $g$（一般级更低），使得对算子 $T(n)$（$(n,q)=1$）有相同的特征值。因此确实可以用本原形式获得关于整个尖点形式空间的信息，必要时使用较低级的本原形式。

此外，人们或许会觉得本原形式的原始定义并不方便，因为它要求 $f\in S_{k}(q,\chi)^{*}$，而后者是“否定地”定义的。然而 W. Li 证明了推论的逆命题：若 $f$ 满足 $f\mid_{k}W=\eta f$ 且 $L(f,s)$ 有 Euler 乘积（它自动也满足所需的函数方程），则 $f$ 是本原形式。这与更大的“自守”哲学完全吻合：“好的”L 函数对应于正确的“本原”对象。

\section{其他的“开场白”}\label{c3:sec:6}

我们像开篇那样结束本章：就这一经典理论中若干进一步的课题作三点简短评注；这些课题对 L 函数与自守形式这一主题的演进影响甚大，并将在其他各章中以（或多或少经过变形的）面貌出现。

\subsection{逆定理}\label{c3:sec:6.1}

导出命题 4.4 证明的那些简单步骤显然可以倒过来，从而给出两个 $1$ 周期全纯函数 $f(z)$ 与 $g(z)$ 由 Fricke 对合相联系（即 $g=f\mid_{k}W$）的判别准则（对 Maass 形式类似，只要 $f,g$ 具有 (4.3) 型的 Fourier 展开）。

对级 $1$，此时 $W\in\mathrm{SL}(2,\mathbb{Z})$，还知道矩阵
\[
\begin{pmatrix}1&1\\0&1\end{pmatrix}\quad\text{与}\quad\begin{pmatrix}0&-1\\1&0\end{pmatrix}
\]
生成 $\mathrm{SL}(2,\mathbb{Z})$；于是由其 L 函数的函数方程，可以推出级 $1$ 模形式的一个\emph{刻画}。

人们希望对更高级得到类似的刻画，但群论不那么简单。Weil 找到了正确的推广：他不仅用 $f$ 本身的 L 函数的函数方程，还用 $f$ 经 Dirichlet 特征扭之后的 L 函数的函数方程。粗略地说，人们不仅需要 $f$ 的 L 级数的信息，还需要 $f$ 的系数 $a_{f}(n)$ 在给定同余类中所生成的 Dirichlet 级数的信息（尽管后者是用 Dirichlet 特征“对偶地”看到的）。

我们讨论全纯情形；非全纯情形类似，只是要用 (4.13) 型的函数方程来“指定”特征值。

例如设 $f\in S_{k}(q,\chi)$，$\psi$ 为模 $m$ 的 Dirichlet 特征。定义
\[
(f\otimes\psi)(z)=\sum_{n\ge1}a_{f}(n)\psi(n)e(nz).
\]
把 $\psi$ 分解为加法特征（这会引入 $\psi$ 的 Gauss 和），可以证明这仍是模形式，确切地说 $f\otimes\psi\in S_{k}(qm^{2},\chi\psi^{2})$。若 $(q,m)=1$，计算 $(f\otimes\psi)\mid_{k}W$ 便得到 $L(f\otimes\psi,s)$ 的如下函数方程，其中 $g=f\mid_{k}W$：
\[
\Lambda(f\otimes\psi,s)=\varepsilon(f\otimes\psi)(qm^{2})^{k/2-s}\Lambda(g\otimes\overline{\psi},k-s),
\]
其中
\begin{equation}\tag{6.1}\label{c3:eq:6.1}
\varepsilon(f\otimes\psi)=i^{k}\chi(m)\psi(q)\tau(\psi)^{2}r^{-1}.
\end{equation}

\begin{annote}
原书此处出现 $r^{-1}$，而全章并未定义 $r$。由函数方程中已显含的因子 $(qm^{2})^{k/2-s}$ 可知 $\varepsilon(f\otimes\psi)$ 必须是模 $1$ 的复数；又 $|\tau(\psi)|=\sqrt{m}$，故 $|\tau(\psi)^{2}|=m$，与之相配的只能是 $m^{-1}$。疑为 $m^{-1}$ 之误，译文照录存疑。
\par\vspace{0.15em}{\raggedleft 【来源：原书 §6.1（A）；Iwaniec《Topics in Classical Automorphic Forms》7.4 节（A）】\par}
\end{annote}
Weil 的定理说：只要对“足够多”的扭都有这样的函数方程，就能保证一个 Dirichlet 级数来自模形式。为简单起见，我们陈述尖点形式的版本。证明见例如 [I1, 7.4]。

\par\medskip\noindent\textbf{定理 6.1}\quad 设
\[
L_{1}(s)=\sum_{n\ge1}a(n)n^{-s}\qquad\text{与}\qquad L_{2}(s)=\sum_{n\ge1}b(n)n^{-s}
\]
是两个对 $\operatorname{Re}(s)>C$（某 $C>0$）绝对收敛的 Dirichlet 级数。

假设存在整数 $k\ge1$、$q\ge1$ 与 $M>0$，使得对任一模 $m$ 的本原特征 $\psi$（$(m,Mq)=1$），Dirichlet 级数
\[
L(f\otimes\psi,s)=\sum_{n\ge1}a(n)\psi(n)n^{-s}\qquad\text{与}\qquad L(g\otimes\psi,s)=\sum_{n\ge1}b(n)\psi(n)n^{-s}
\]
均可解析延拓为竖直带内有界的整函数，使得
\[
\Lambda(f\otimes\psi,s)=(2\pi)^{-s}\Gamma(s)L(f\otimes\psi,s)
\]
与
\[
\Lambda(g\otimes\psi,s)=(2\pi)^{-s}\Gamma(s)L(g\otimes\psi,s)
\]
都是整函数，并满足函数方程
\[
\Lambda(f\otimes\psi,s)=\varepsilon(f\otimes\psi)(qm^{2})^{k/2-s}\Lambda(g\otimes\overline{\psi},k-s),
\]
其中 $\varepsilon(f\otimes\psi)$ 由 (6.1) 给出。

则存在模某 $q$ 的特征 $\chi$ 的尖点形式 $f\in S_{k}(q,\chi)$，使得 $L(f,s)=L_{1}(s)$ 且 $L(f\mid_{k}W,s)=L_{2}(s)$。

这类逆定理的推广有很长的历史，并对自守形式的许多发展、特别是 Langlands 的函子性猜想极为重要；可以举出 Gelbart 与 Jacquet 关于 $\mathrm{GL}(2)$ 形式的对称平方的构造。Cogdell 与 Piatetski-Shapiro 较新的结果，也被用于 Lafforgue 关于函数域上 $\mathrm{GL}(n)$ 的整体 Langlands 对应的证明。

\subsection{Rankin--Selberg L 函数}\label{c3:sec:6.2}

设 $f$ 是级为 $q$ 的权 $k$ 尖点形式。我们在 (4.8) 中提到，上界
\[
\sum_{n\le N}|a_{f}(n)|^{2}\ll N^{k}\qquad(N\ge1)
\]
相当容易得到。Rankin 与 Selberg 各自独立地证明了它确实是精确的；如同 Euler 乘积 (5.1) 的证明一样，他们的方法对后来 L 函数与自守形式理论的发展有相当大的影响。

给定两个模形式 $f$ 与 $g$（权与级任意，其中至多一个可以是非全纯的），先定义（朴素的）Rankin--Selberg L 函数
\[
L(f\times g,s)=\sum_{n\ge1}a_{f}(n)a_{g}(n)n^{-s},
\]
于是 $L(f\times\bar f,s)$ 的系数就是 $|a_{f}(n)|^{2}$。

由 Fourier 系数的多项式界，该级数在 $\operatorname{Re}(s)$ 充分大时收敛。Rankin 与 Selberg 工作的要点在于：这类新的 L 函数具有 Hecke L 函数（交换情形）与自守 L 函数的大部分良好解析性质；即它们有解析延拓、函数方程，并在有利情形下有 Euler 乘积展开。这里我们只陈述最简单的情形：

\par\medskip\noindent\textbf{命题 6.2}\quad 设 $f\in S_{k}(1)$，$g\in S_{\ell}(1)$，$k\ge2$，$\ell\ge2$，为两个全纯模形式。除 $L(f\times g,s)$ 外，再定义
\[
L(f\otimes g,s)=\zeta\bigl(2s-(k+\ell)+2\bigr)L(f\times g,s)
\]
与
\[
\Lambda(f\otimes g,s)=(2\pi)^{-2s}\Gamma(s)\Gamma\bigl(s-(k+\ell)/2+1\bigr)L(f\otimes g,s).
\]
则 $\Lambda(f\otimes g,s)$ 可解析延拓到 $\mathbb{C}$；若 $f\ne\bar g$ 则它是整函数，否则在 $s=0$ 与 $s=k+\ell$ 处有单极点，并满足函数方程
\[
\Lambda(f\otimes g,s)=\Lambda(\bar f\otimes\bar g,k+\ell-s).
\]

于是，若 $f\ne\bar g$，则 $L(f\otimes g,s)$ 是整函数；若 $f=\bar g$，则它在 $s=k=\ell$ 处有单极点。留数由下式给出：
\[
\operatorname{Res}_{s=k}L(f\times g,s)=\frac{3}{\pi}\frac{(4\pi)^{k}}{\Gamma(k)}\langle f,f\rangle.
\]
\emph{证明梗概}\quad 证明关键地依赖于 (3.3) 中定义的非全纯 Eisenstein 级数 $E(z;s)$；确切地说，关键在于 Rankin--Selberg 积分表示
\begin{align*}
&(4\pi)^{-s}\Gamma(s)L(f\times g,s)\\
&\quad=\int_{\Gamma_{0}(1)\backslash\mathbf{H}}f(z)\overline{g(z)}y^{(k+\ell)/2}E\bigl(z;s-(k+\ell)/2+1\bigr)\,d\mu(z),
\end{align*}
其中出现非全纯 Eisenstein 级数。为证明它，注意
\[
(4\pi)^{-s}\Gamma(s)n^{-s}=\int_{0}^{+\infty}e^{-4\pi ny}y^{s}\,\frac{dy}{y},
\]
以及
\[
\int_{0}^{1}f(z)\overline{g(z)}\,dx=\sum_{n\ge1}a_{f}(n)\overline{a_{g}(n)}e^{-4\pi ny};
\]
于是，由于 $]0,1[\times]0,+\infty[\subset\mathbb{R}^{2}$ 是 $\mathrm{SL}(2,\mathbb{Z})$ 中 $+\infty$ 的稳定化子 $\Gamma_{\infty}$ 的基本域，我们得到
\begin{align*}
&(4\pi)^{-s}\Gamma(s)L(f\times g,s)=\int_{\Gamma_{\infty}\backslash\mathbf{H}}f(z)\overline{g(z)}y^{s-1}\,dx\,dy\\
&\quad=\sum_{\gamma\in\Gamma_{\infty}\backslash\mathrm{SL}(2,\mathbb{Z})}\int_{\mathrm{SL}(2,\mathbb{Z})\backslash\mathbf{H}}f(z)\overline{g(z)}y^{(k+\ell)/2}(\operatorname{Im}\gamma z)^{s-(k+\ell)/2+1}\,d\mu(z),
\end{align*}
这里用到 $z\mapsto f(z)y^{k/2}$ 与 $z\mapsto g(z)y^{\ell/2}$ 都是 $\mathrm{SL}(2,\mathbb{Z})$ 不变函数这一事实。交换求和次序，由 $E(z;s)$ 的定义即得上述积分公式。

再把两边乘以 $\pi^{-s}\Gamma(s-(k+\ell)/2+1)\zeta(2s-(k+\ell)+2)$，左边便得到 $\Lambda(f\otimes g,s)$。右边由 $E(z;s)$ 的函数方程 (4.7) 可知，我们得到一个在 $s\mapsto(k+\ell)/2-s$ 下不变的表达式，于是即得 $\Lambda(f\otimes g,s)$ 的函数方程。

\begin{annote}
原文说该表达式在 $s\mapsto(k+\ell)/2-s$ 下不变，这与命题 6.2 的函数方程 $\Lambda(f\otimes g,s)=\Lambda(\bar f\otimes\bar g,k+\ell-s)$ 不符：$(k+\ell)/2$ 是临界线的位置，不是对称轴，对称轴应为 $(k+\ell)/2$ 所对应的 $s\mapsto k+\ell-s$。译文照录原书，读者按后者理解即可。
\par\vspace{0.15em}{\raggedleft 【来源：原书 §6.2（A）】\par}
\end{annote}

由于 $E(z,s)$ 在 $s=1$ 处有单极点，留数为 $1/\operatorname{Vol}(\mathrm{SL}(2,\mathbb{Z})\backslash\mathbf{H})=3/\pi$，其余论断也都是这一积分表示的推论。

特别要注意，具有“良好”解析性质的确实是 $L(f\otimes g,s)$。较简单的 $L(f\times g,s)$ 固然在 $\mathbb{C}$ 上亚纯，但它在 $s=\rho/2+(k+\ell)/2+1$（$\rho$ 为 Riemann zeta 函数的零点）处有无穷多个极点。\hfill$\square$

\par\medskip\noindent\textbf{推论 6.3}\quad 设 $f\in S_{k}(1)$ 为模形式。则当 $X\to+\infty$ 时
\[
\sum_{n\le X}|a_{f}(n)|^{2}\sim c(f)X^{k},
\]
其中
\[
c(f)=\operatorname{Res}_{s=k}L(f\times\bar f,s).
\]
\emph{证明}\quad 对 $L(f\times\bar f,s)$ 施以标准的围道积分方法：$s=k$ 处的单极点给出主项。\hfill$\square$

\par\medskip\noindent\textbf{注 6.4}\quad 除上述之外，Selberg 还注意到：若 $f$ 与 $g$ 的 Fourier 系数是乘性的，则 $L(f\times g,s)$ 的系数也是乘性的，故这个 L 函数也有 Euler 乘积展开。设 $f\in S_{k}(1)$ 与 $g\in S_{\ell}(1)$ 为本原形式，利用 (5.4) 写作
\[
a_{f}(p^{k})=\alpha_{1}^{k}+\alpha_{2}^{k},\qquad a_{g}(p^{k})=\beta_{1}^{k}+\beta_{2}^{k},
\]
其中 $(\alpha_{1},\alpha_{2})$（相应地 $(\beta_{1},\beta_{2})$）是二次方程
\[
1-a_{f}(p)X+p^{k-1}X^{2}=0\quad(\text{相应地 }1-a_{g}(p)X+p^{\ell-1}X^{2}=0)
\]
的根。则
\begin{align*}
\sum_{k\ge0}a_{f}(p^{k})a_{g}(p^{k})X^{k}&=\frac{1}{1-\alpha_{1}\beta_{1}X}+\frac{1}{1-\alpha_{1}\beta_{2}X}+\frac{1}{1-\alpha_{2}\beta_{1}X}+\frac{1}{1-\alpha_{2}\beta_{2}X}\\
&=\frac{1-p^{\ell+k-2}X^{2}}{(1-\alpha_{1}\beta_{1}X)(1-\alpha_{1}\beta_{2}X)(1-\alpha_{2}\beta_{1}X)(1-\alpha_{2}\beta_{2}X)};
\end{align*}
于是当 $\operatorname{Re}(s)$ 充分大时，$L(f\otimes g,s)$ 有如下 Euler 乘积展开：
\begin{align*}
L(f\otimes g,s)=\prod_{p}\Bigl(&(1-\alpha_{1,p}\beta_{1,p}p^{-s})(1-\alpha_{1,p}\beta_{2,p}p^{-s})\\
&\cdot(1-\alpha_{2,p}\beta_{1,p}p^{-s})(1-\alpha_{2,p}\beta_{2,p}p^{-s})\Bigr)^{-1}.
\end{align*}

一般情形下，若 $f$ 与 $g$ 都是本原形式，同样的形式计算表明 $L(f\otimes g,s)$ 有 Euler 乘积展开，且除有限多个因子外所有 Euler 因子都具上述形状。不过若走古典路线，函数方程就变得相当棘手；因为它应当从 Eisenstein 级数继承而来，而对 Eisenstein 级数，只要所涉及的群有多于一个尖点，函数方程就取矩阵形式，需要用到所有尖点处的 Eisenstein 级数。用阿代尔语言，结果又要干净得多，证明在概念上大概也更简单。

\par\medskip\noindent\textbf{注 6.5}\quad 为与交换 L 函数的类比保持一致，应当把 $L(f\otimes g,s)$ 看作 $L(\chi_{1}\chi_{2},s)$（$\chi_{i}$ 为某些 Dirichlet 特征）的类比。当然，后一情形仍是交换 L 函数，因为 $1\times1=1$。特别地，应当把 $L(f\otimes g,s)$ 看作衡量两个模形式之间“距离”的一种方式。$s=1$ 处极点的阶数给出了判定 $f=\bar g$ 与否的办法，这非常有用；$s=k$ 处的留数与 $f$ 的 Petersson 范数相关，这一点同样有用（例如本原形式的 Petersson 范数往往正是这样研究最好）。

\begin{annote}
原文作“$s=1$ 处极点的阶数”，应为 $s=k$。$L(f\times g,s)$ 的极点来自 $E(z;s)$ 在 $s=1$ 处的极点，经积分表示中的平移 $s\mapsto s-(k+\ell)/2+1$ 后落在 $s=k$（取 $k=\ell$）处；命题 6.2 及其后的留数公式也都指向 $s=k$。译文照录原书。
\par\vspace{0.15em}{\raggedleft 【来源：原书 §6.2（A）】\par}
\end{annote}

一个差别是记法 $f\otimes g$ 颇为形式化。直到最近，Ramakrishnan 才证明 Rankin--Selberg L 函数也附着于某个自守对象，即 $\mathrm{GL}(2)\times\mathrm{GL}(2)$ 上的一个自守形式，正如 Langlands 函子性猜想所预言的那样。

\subsection{重访 theta 函数与二次域}\label{c3:sec:6.3}

前面已提到，用来证明交换 L 函数解析性质的那类 theta 函数本身就是模形式。Hecke 以及后来的 Maass 考虑了一个特别有趣的特例，它可以看作把二次域 $K$ 上 $\mathrm{GL}(1)$ 的自守形式与 $\mathbb{Q}$ 上 $\mathrm{GL}(2)$ 的自守形式相联系的第一例。在类群特征（即权平凡、模 $\mathfrak{m}=\mathcal{O}$）这一特殊情形下，结果可陈述如下：

\par\medskip\noindent\textbf{命题 6.6}\quad 设 $K/\mathbb{Q}$ 是判别式为 $D$ 的二次扩张。设 $\chi$ 是 $K$ 的理想类群的特征。则存在级为 $|D|$、nebentypus 为 $\varepsilon_{D}$（等于 $K$ 相伴的二次特征）的本原模形式 $f_{\chi}$；当 $D<0$ 时它是权 $1$ 的全纯形式，当 $D>0$ 时它是特征值 $\lambda=1/4$ 的 Maass 形式，使得
\[
L(f_{\chi},s)=L(\chi,s),
\]
且这即使在每个 $p$ 因子上也局部成立。此外，$f_{\chi}$ 是尖点形式当且仅当特征 $\chi$ 不是实特征。

可以用带特征的 theta 函数构造 $f$：
\[
f(z)=\sum_{\mathfrak{a}\ne0}\chi(\mathfrak{a})e(zN\mathfrak{a})
\]
（当 $\chi$ 非实且 $K$ 为虚二次域时）。

或者，也容易看出 $\chi$ 的 L 函数的函数方程确实与 $f$ 的存在相容。展开
\[
L(\chi,s)=\sum_{n\ge1}a_{\chi}(n)n^{-s},
\]
便可应用 Weil 的逆定理：$K$ 上交换 L 函数的函数方程表明该定理的假设得到满足。$f_{\chi}$ 的本原性容易推出，因为按构造 $L(f_{\chi},s)$ 有 Euler 乘积。

在实域情形，这基本上给出了已知的 Maass 形式的唯一显式例子。注意它们的 Fourier 系数是代数数。人们猜想：所有特征值 $\lambda=1/4$ 的本原 Maass 形式都必来自实二次域。

类似的论证可用于其他权的特征，特别地给出数域上具有复乘的椭圆曲线的 Hasse--Weil zeta 函数的解析延拓与函数方程（Deuring、Weil；一个完全详尽的特殊情形见例如 [I2, 8]）。

\section*{参考文献}
\addcontentsline{toc}{section}{参考文献}
\markright{参考文献}

\begin{itemize}[leftmargin=4.6em,itemsep=0.22em,topsep=0.3em,label={}]
\item[{[AL]}] A. Atkin and J. Lehner, Hecke operators on $\Gamma_{0}(m)$, \emph{Math.\ Ann.}, \textbf{185}(1970), 134--160.
\item[{[Bu]}] D. Bump, \emph{Automorphic Forms and Representations}, Cambridge University Press, Cambridge, UK, 1996.
\item[{[CF]}] J. W. S. Cassels and A. Fröhlich, eds., \emph{Algebraic Number Theory}, Academic Press, New York, 1967.
\item[{[CP]}] J. Cogdell and I. Piatetski-Shapiro, \emph{The Arithmetic and Spectral Analysis of Poincaré Series}, Perspectives in Mathematics 13, Academic Press, New York, 1990.
\item[{[Co]}] H. Cohen, \emph{Advanced Topics in Computational Number Theory}, Graduate Texts in Mathematics 193, Springer-Verlag, New York, 2000.
\item[{[CS]}] B. Conrey and K. Soundararajan, Real zeros of quadratic Dirichlet $L$-functions, \emph{Invent.\ Math.}, \textbf{150}-1(2002), 1--44.
\item[{[De]}] P. Deligne, Les conjectures de Weil I, \emph{Inst.\ Hautes Études Sci.\ Publ.\ Math.}, \textbf{43}(1974), 273--300.
\item[{[Di]}] P. G. Dirichlet, Beweis des Satzes, dass jede unbegrenzte arithmetische Progression, deren erstes Glied und Differenz ganze Zahlen ohne gemeinschaftlichen Factor sind, unendlich viele Primzahlen enthält, in \emph{Collected Works}, Vol.~I, Chelsea, New York, 1969, 313--342.
\item[{[Du]}] W. Duke, The dimension of the space of cusp forms of weight one, \emph{Internat.\ Math.\ Res.\ Notices}, \textbf{2}(1995), 99--109.
\item[{[Ge]}] S. Gelbart, \emph{Automorphic Forms on Adèle Groups}, Annals of Mathematical Studies 83, Princeton University Press, Princeton, NJ, 1975.
\item[{[Go]}] D. Goldfeld, The class number of quadratic fields and the conjectures of Birch and Swinnerton-Dyer, \emph{Ann.\ Scuola Norm.\ Sup.\ Pisa} 3, \textbf{4}(1976), 623--663.

\item[{[I1]}] H. Iwaniec, \emph{Topics in Classical Automorphic Forms}, Graduate Studies in Mathematics 17, American Mathematical Society, Providence, 1997.
\item[{[I2]}] H. Iwaniec, \emph{Introduction to the Spectral Theory of Automorphic Forms}, Biblioteca de la Revista Matemática Iberoamericana, Madrid, 1995.
\item[{[IR]}] K. Ireland and M. Rosen, \emph{A Classical Introduction to Modern Number Theory}, 2nd ed., Graduate Texts in Mathematics 84, Springer-Verlag, New York, 1990.
\item[{[KS]}] N. Katz and P. Sarnak, \emph{Random Matrices, Frobenius Eigenvalues, and Monodromy}, Colloquium Publications, American Mathematical Society, Providence, 1998.
\item[{[L]}] E. Landau, Bemerkungen zum Heilbronnschen Satz, \emph{Acta Arithmetica}, \textbf{1}(1936), 1--18.
\item[{[La]}] S. Lang, \emph{Algebraic Number Theory}, Graduate Texts in Mathematics 110, Springer-Verlag, New York, 1986.
\item[{[LRS]}] W. Luo, Z. Rudnick, and P. Sarnak, On Selberg's eigenvalue conjecture, \emph{Geometric Functional Anal.}, \textbf{5}(1994), 387--401.
\item[{[Mi]}] T. Miyake, \emph{Modular Forms}, Springer-Verlag, New York, 1989.
\item[{[PS]}] R. Phillips and P. Sarnak, On cusp forms for cofinite subgroups of $PSL(2,\mathbf{R})$, \emph{Invent.\ Math.}, \textbf{80}(1985), 339--364.
\item[{[Rie]}] B. Riemann, Über die Anzahl der Primzahlen unter einer gegebenen Grösse, in \emph{Monatsberichte der Berliner Akademie}, Berliner Akademie, Berlin, 1859, 671--680; reprinted in R. Narasimhan, ed., \emph{Gesammelte mathematische Werke, wissenschaftlicher Nachlass und Nachträge}, Springer-Verlag, Berlin, 1990, 177--185.
\item[{[RS]}] Z. Rudnick and P. Sarnak, Zeros of principal $L$-functions and random matrix theory, \emph{Duke Math.\ J.}, \textbf{81}(1996), 269--322.
\item[{[Sa]}] P. Sarnak, Quantum chaos, symmetry and zeta functions, \emph{Curr.\ Dev.\ Math.}, \textbf{3}(1997), 84--115.
\item[{[Se1]}] J.-P. Serre, \emph{Cours d'arithmétique}, 3rd ed., Presses Université de France, Paris, 1988.
\item[{[Se2]}] J.-P. Serre, Modular forms of weight one and Galois representations, in A. Fröhlich, ed., \emph{Algebraic Number Fields}, Academic Press, New York, 1977, 193--268.
\item[{[Sh]}] G. Shimura, \emph{Introduction to the Arithmetic Theory of Automorphic Functions}, Iwanami Shoten and Princeton University Press, Princeton, NJ, 1971.

\item[{[Si]}] C. L. Siegel, Über die Classenzahl quadratischer Zahlkörper, \emph{Acta Arithmetica}, \textbf{1}(1936), 83--86.
\item[{[Ta]}] J. Tate, Fourier analysis in number fields and Hecke's zeta function, in J. W. S. Cassels and A. Fröhlich, eds., \emph{Algebraic Number Theory}, Academic Press, New York, 1967, 305--347 [CF].
\item[{[Tu]}] J. Tunnell, Graduate course, Rutgers University, New Brunswick, NJ, 1995--1996.
\item[{[Wa]}] L. Washington, \emph{Cyclotomic Fields}, 2nd ed., Graduate Texts in Mathematics 83, Springer-Verlag, New York, 1997.
\item[{[We]}] A. Weil, Prehistory of the zeta function, in K. E. Aubert, E. Bombieri, and D. Goldfeld, eds., \emph{Number Theory, Trace Formulas and Discrete Groups}, Academic Press, New York, 1989.
\item[{[Ti]}] A. C. Titchmarsh, \emph{The Theory of the Riemann Zeta-Function}, 2nd ed. (revised by D. R. Heath-Brown), Oxford University Press, Oxford, UK, 1986.
\end{itemize}

